\documentclass[11pt,a4paper]{article}

% ============================================================
% Uniform MTDF preamble -- identical across all papers
% ============================================================
\usepackage[utf8]{inputenc}
\usepackage[T1]{fontenc}
\usepackage{lmodern}
\usepackage[english]{babel}
\usepackage[a4paper,margin=2.2cm]{geometry}
\usepackage{amsmath,amssymb,amsthm}
\usepackage{graphicx}
\usepackage{float}
\usepackage{booktabs}
\usepackage{array}
\usepackage{longtable}
\usepackage{tabularx}
\usepackage{multirow}
\usepackage{caption}
\usepackage[dvipsnames,table]{xcolor}
\usepackage{mdframed}
\usepackage{parskip}
\usepackage{ragged2e}
\usepackage[hidelinks,breaklinks=true]{hyperref}
\usepackage{microtype}

% Colours matching the HTML theme
\definecolor{accent}{HTML}{1A4B8A}
\definecolor{callout-bg}{HTML}{FBFBFB}
\definecolor{tablehead}{HTML}{F0F0F0}

% Prevent any table from overflowing
\setlength{\tabcolsep}{4pt}
\renewcommand{\arraystretch}{1.15}

% Caption style (compact, italic, small like the HTML)
\captionsetup{font=small,labelfont=bf,labelsep=period,justification=centerfirst}

% Hyperref metadata placeholders (filled per-paper below)
\hypersetup{
  pdftitle={MTDF Independent GPU Validation: Phases 1-6},
  pdfauthor={Ingo Mesche},
  pdfsubject={MTDF - Mesche's Tensor Dynamics Framework},
  colorlinks=false
}

% Prevent overfull hbox from long URLs/DOIs in bibliography
\sloppy
\emergencystretch=3em

% Tolerant line breaking so long identifiers (Python filenames etc.) don't
% produce large overfull hboxes in narrow table columns. These settings keep
% all content on-page while slightly relaxing right-margin strictness.
\tolerance=1200
\hbadness=2000
\hyphenpenalty=150
\exhyphenpenalty=50

% ============================================================
\title{MTDF Independent GPU Validation: Phases 1--6}
\author{Ingo Mesche\\ \small Independent Researcher, Malta \\ \small \texttt{imesche@outlook.com}}
\date{V74 / February 2026}

\begin{document}
\maketitle
\begin{center}\itshape Blind reproduction, CMB consistency, environment signal confirmation, sensitivity forecasts, full Planck hard-falsifier test, and cross-probe discriminators\end{center}

\vspace{0.4em}

\begin{center}\itshape Hardware: NVIDIA RTX 4080 SUPER (16 GB VRAM, CUDA 12.9)\end{center}



\begin{mdframed}[linecolor=accent,linewidth=0.5pt,backgroundcolor=callout-bg,leftline=true,rightline=true,topline=true,bottomline=true,innerleftmargin=8pt,innerrightmargin=8pt,innertopmargin=6pt,innerbottommargin=6pt]
\textbf{Abstract.}
This document reports the results of a fully independent, GPU-accelerated validation of Mesche's Tensor Dynamics Framework (MTDF). Six phases were executed over multiple sessions using code written from scratch, with no shared libraries or routines from the original MTDF analysis pipeline. Phase~1 reproduces all V18 dashboard $\chi$$^{2}$ values to four decimal places. Phase~2 tests MTDF against the Planck 2018 CMB power spectrum via MCMC with a perturbative correction layer, finding $\Delta$$\chi$$^{2}$~<~1 and an unconstrained early-field-energy amplitude. Phase~3 independently confirms the SN~$\times$~void environment signal reported in MTDF\_\allowbreak{}02, with bootstrap confidence intervals excluding zero. Phase~4 forecasts that 5$\sigma$ detection of the environment signal requires approximately 3,400 low-redshift supernovae, placing definitive confirmation within the reach of LSST/\allowbreak{}Rubin. Phase~5 executes the Priority~1 hard falsifier: the full Planck 2018 plik TTTEEE + lowl + lensing likelihood using the \texttt{class\_\allowbreak{}mtdf} Boltzmann solver with modified perturbation equations, finding $\Delta$$\chi$$^{2}$~=~+0.63. This MTDF implementation is statistically indistinguishable from $\Lambda$CDM at the CMB within the precision and assumptions of this likelihood comparison. Phase~6 extends validation to cross-probe discriminators: a redshift transition scan (3.6$\sigma$ at z~<~0.03), weak lensing $\times$ environment (KiDS-1000; systematics-clean null with explicit upper limits), derived parameter consistency (S\textsubscript{8} tension reduced), and CMB lensing $\times$ voids (Planck~PR4 stacking on DESI~BGS and BOSS~DR12 void catalogues; pipeline validated against published detections).
\end{mdframed}

\tableofcontents

\begin{mdframed}[linecolor=accent,linewidth=0.5pt,backgroundcolor=callout-bg,leftline=true,rightline=true,topline=true,bottomline=true,innerleftmargin=8pt,innerrightmargin=8pt,innertopmargin=6pt,innerbottommargin=6pt]
\textbf{Claim status.}
This paper distinguishes between:
(i) calibration anchors,
(ii) validation targets,
(iii) diagnostic consistency checks, and
(iv) speculative extensions.
Only results explicitly labelled as validation targets are used as evidence for the core MTDF claim. Diagnostic and speculative sections are included to define future tests, not to claim established confirmation.
\end{mdframed}

\section{Scope and methodology}
\label{sec:sec1}
The purpose of this validation exercise is to subject MTDF to a battery of independent tests that are entirely separate from the original analysis pipeline. All code was written from scratch across multiple sessions. No libraries, routines, or intermediate results from the original MTDF codebase were used for Phases~1, 3, 4, or 6. Phase~2 used the CosmoPower $\Lambda$CDM emulator with a bespoke MTDF correction layer. Phase~5 used the existing \texttt{class\_\allowbreak{}mtdf} modified Boltzmann solver through cobaya with the full Planck likelihood stack.

The six phases test qualitatively different aspects of the framework:

\begin{itemize}
  \item \textbf{Phase 1} tests internal consistency: can the published $\chi$$^{2}$ values be reproduced from raw data using independent code?
  \item \textbf{Phase 2} tests CMB compatibility (perturbative): is MTDF allowed by a compressed CMB likelihood?
  \item \textbf{Phase 3} tests the primary empirical prediction: is the SN~$\times$~void environment signal real and reproducible?
  \item \textbf{Phase 4} tests future viability: what survey specifications are needed for definitive detection?
  \item \textbf{Phase 5} tests the hard falsifier: does MTDF survive the full Planck likelihood with modified Boltzmann equations?
  \item \textbf{Phase 6} tests cross-probe discriminators: does the transition signature survive in weak lensing, CMB lensing, and derived parameters?
\end{itemize}

\begin{mdframed}[linecolor=accent,linewidth=0.5pt,backgroundcolor=callout-bg,leftline=true,rightline=true,topline=true,bottomline=true,innerleftmargin=8pt,innerrightmargin=8pt,innertopmargin=6pt,innerbottommargin=6pt]
\textbf{Independence guarantee.}
Phases~1, 3, and 6 use fully independent implementations. Phase~2 uses a perturbative correction layer on top of a public $\Lambda$CDM emulator (CosmoPower), parameterised by a continuous coupling \( k_f \) where \( k_f = 0 \) is $\Lambda$CDM and \( k_f = 1 \) is full MTDF. Phase~4 builds on Phase~3 infrastructure with additional Monte Carlo machinery. Phase~5 uses the existing \texttt{class\_\allowbreak{}mtdf} Boltzmann solver with cobaya and the full Planck plik likelihood. This is the only phase that uses pre-existing MTDF code, which is appropriate since the test is specifically to validate that code against Planck. Phase~6 uses public data products (KiDS-1000 shear catalogues, Planck~PR4 lensing maps, DESI/\allowbreak{}BOSS void catalogues) with all stacking and analysis code written from scratch. All GPU computations were cross-checked against CPU results.
\end{mdframed}

\subsection{Hardware and software}
All computations were performed on a desktop workstation equipped with an NVIDIA RTX 4080 SUPER GPU (16 GB VRAM, CUDA 12.9, compute capability 8.9). The software environment consists of Python 3.12 with cobaya 3.6, emcee, getdist, CuPy, JAX, TensorFlow, and CosmoPower. All datasets were pre-downloaded: Pantheon+ with full covariance matrix, DESI Y1 BAO, cosmic chronometers, DR16 f$\sigma$\textsubscript{8} growth measurements, and DESI void catalogues (VoidFinder, REVOLVER, VIDE).

\section{Phase 1: Independent $\chi$$^{2}$ reproduction}
\label{sec:sec2}
PASSED: All V18 dashboard values reproduced to $\delta$ = 0.000049A fully independent implementation reproduced all V18 dashboard $\chi$$^{2}$ values. The code reads raw data files, constructs likelihood functions from scratch, and evaluates MTDF predictions at the V18 parameter point. No dashboard code was consulted or reused.

\subsection{Strict total}
\begin{table}[H]
\centering
\footnotesize
\begin{tabularx}{\linewidth}{l>{\RaggedRight\arraybackslash}X>{\RaggedRight\arraybackslash}X>{\RaggedRight\arraybackslash}X}
\toprule
\textbf{Metric} & \textbf{Target (V18)} & \textbf{Reproduced} & \textbf{$\delta$} \tabularnewline
$\chi$$^{2}$/$\nu$ & 1.1683 & 1.1683 & 0.000049 \tabularnewline
Degrees of freedom $\nu$ & 1745 & 1745 & -- \tabularnewline
\bottomrule
\end{tabularx}
\end{table}

\subsection{Vector pillar breakdown}
\begin{table}[H]
\centering
\footnotesize
\begin{tabularx}{\linewidth}{l>{\RaggedRight\arraybackslash}X>{\RaggedRight\arraybackslash}X>{\RaggedRight\arraybackslash}X}
\toprule
\textbf{Pillar} & \textbf{$\chi$$^{2}$} & \textbf{DOF} & \textbf{$\chi$$^{2}$/$\nu$} \tabularnewline
Pantheon+ SNe Ia & 2012.72 & 1700 & 1.184 \tabularnewline
DESI Y1 BAO & 15.39 & 12 & 1.282 \tabularnewline
Cosmic chronometers H(z) & 6.98 & 15 & 0.466 \tabularnewline
DR16 f$\sigma$\textsubscript{8} growth & 2.00 & 3 & 0.667 \tabularnewline
CMB distance prior (diagnostic only) & 595.33 & 3 & 198.44 \tabularnewline
\bottomrule
\end{tabularx}
\end{table}

Plus 15 scalar pillars contributing $\chi$$^{2}$ = 1.67, matching exactly.

\begin{mdframed}[linecolor=accent,linewidth=0.5pt,backgroundcolor=callout-bg,leftline=true,rightline=true,topline=true,bottomline=true,innerleftmargin=8pt,innerrightmargin=8pt,innertopmargin=6pt,innerbottommargin=6pt]
\textbf{Significance.} Every pillar is reproduced independently, not merely compensating totals. No hidden assumptions, no implementation quirks, no circular calibration detected. The V18 strict protocol is validated by independent audit.
\end{mdframed}

\section{Phase 2: Planck CMB consistency}
\label{sec:sec3}
PASSED: $\Delta$$\chi$$^{2}$ < 1; \, k\textsubscript{f} = 1 within 95\% CI in all runsPhase 2 tests whether MTDF is compatible with the Planck 2018 CMB power spectrum. The approach uses the CosmoPower neural network emulator for $\Lambda$CDM C\textsubscript{$\ell$} spectra, augmented by an analytic MTDF correction layer that modifies the sound horizon, angular diameter distance, and ISW contribution as functions of a continuous coupling parameter \( k_f \). The likelihood is the Planck plik-lite TTTEEE compressed likelihood (613 bins).

\subsection{Emulator validation}
\begin{table}[H]
\centering
\footnotesize
\begin{tabularx}{\linewidth}{l>{\RaggedRight\arraybackslash}X>{\RaggedRight\arraybackslash}X>{\RaggedRight\arraybackslash}X}
\toprule
\textbf{Spectrum} & \textbf{Prediction time} & \textbf{Mean error vs CAMB} & \textbf{Max error} \tabularnewline
TT & 7.5 ms & 0.09\% & 0.20\% \tabularnewline
TE & 4.1 ms & 0.46\% & 15.4\%* \tabularnewline
EE & 7.6 ms & 0.11\% & 0.29\% \tabularnewline
\bottomrule
\end{tabularx}
\end{table}

*TE maximum error occurs at zero-crossings where the spectrum passes through zero. This is expected and numerically harmless.

\subsection{Fixed-parameter comparison}
\begin{table}[H]
\centering
\footnotesize
\begin{tabularx}{\linewidth}{l>{\RaggedRight\arraybackslash}X>{\RaggedRight\arraybackslash}X>{\RaggedRight\arraybackslash}X}
\toprule
\textbf{Cosmology} & \textbf{$\chi$$^{2}$} & \textbf{Bins} & \textbf{$\chi$$^{2}$/bin} \tabularnewline
$\Lambda$CDM (CAMB reference) & 584.45 & 613 & 0.953 \tabularnewline
MTDF at k\textsubscript{f} = 1 (corrected) & 585.26 & 613 & 0.955 \tabularnewline
$\Delta$$\chi$$^{2}$ & +0.81 &  &  \tabularnewline
\bottomrule
\end{tabularx}
\end{table}

$\Lambda$CDM $\chi$$^{2}$ = 584.45 matches the Planck 2018 published value ($\sim$583). Pipeline correctly wired.

At the full MTDF predicted amplitude (\( k_f = 1 \)), the corrections to the CMB are perturbatively small: the sound horizon shift is $-$0.071\%, the angular scale shift is +0.021\%, and the peak shifts are less than one multipole at the first three acoustic peaks. The $\Delta$$\chi$$^{2}$ of +0.81 is far below any statistical significance threshold.

\subsection{Normalisation correction}
During development, a normalisation inconsistency was identified and corrected. Three different \( k_f \) conventions existed in the codebase (see Appendix). After correction, the physical mapping is:

\[
f_{\rm kick} = \frac{\lambda_{\rm MTDF}}{24} = \frac{(1 - \beta_{\rm eos})^2}{(1+\alpha)\times 24} = 0.003303
\]

At \( k_f = 1 \): $\theta$\textsubscript{s} ratio = 1.000212 (shift of +0.021\%), ISW boost = 0.25\%. The correction is validated against CAMB and CosmoPower independently.

\subsection{Production MCMC results}
Three production MCMC runs were completed (5000 steps, 32 walkers each), using the emcee affine-invariant sampler with GPU-accelerated likelihood evaluation.

\subsubsection{Model comparison}
\begin{table}[H]
\centering
\footnotesize
\begin{tabularx}{\linewidth}{l>{\RaggedRight\arraybackslash}X>{\RaggedRight\arraybackslash}X>{\RaggedRight\arraybackslash}X}
\toprule
\textbf{Metric} & \textbf{Planck-only} & \textbf{Planck + BAO + SN} & \textbf{Planck + BAO + SN (k\textsubscript{f}\,<\,2)} \tabularnewline
N\textsubscript{data} & 613 & 2,326 & 2,326 \tabularnewline
$\Lambda$CDM $\chi$$^{2}$ & 583.10 & 2,607.80 & 2,607.80 \tabularnewline
MTDF best-fit $\chi$$^{2}$ & 582.24 & 2,607.10 & 2,606.89 \tabularnewline
$\Delta$$\chi$$^{2}$ & $-$0.85 & $-$0.70 & $-$0.91 \tabularnewline
$\Delta$AIC & +1.15 & +1.30 & +1.09 \tabularnewline
\bottomrule
\end{tabularx}
\end{table}

$\Delta$$\chi$$^{2}$ negative = MTDF marginally better in raw fit. $\Delta$AIC positive = $\Lambda$CDM mildly preferred by Occam penalty for the extra k\textsubscript{f} parameter. Result: completely inconclusive; neither model is favoured.

\subsubsection{k\textsubscript{f} posterior}
\begin{table}[H]
\centering
\footnotesize
\begin{tabularx}{\linewidth}{l>{\RaggedRight\arraybackslash}X>{\RaggedRight\arraybackslash}X>{\RaggedRight\arraybackslash}X}
\toprule
\textbf{} & \textbf{Planck-only [0,\,5]} & \textbf{Combined [0,\,5]} & \textbf{Combined [0,\,2]} \tabularnewline
Best-fit & 4.63 & 0.085 & 0.086 \tabularnewline
Median & 2.63 & 1.93 & \textbf{0.93} \tabularnewline
68\% CI & {}[0.89, 4.28] & {}[0.59, 3.72] & {}[0.28, 1.65] \tabularnewline
95\% UL & 4.77 & 4.53 & 1.89 \tabularnewline
k\textsubscript{f}\,=\,1 in 95\%? & Yes & Yes & Yes \tabularnewline
k\textsubscript{f}\,=\,0 in 68\%? & Yes & Yes & Yes \tabularnewline
\bottomrule
\end{tabularx}
\end{table}

\subsubsection{H\textsubscript{0} posterior}
\begin{table}[H]
\centering
\footnotesize
\begin{tabularx}{\linewidth}{l>{\RaggedRight\arraybackslash}X>{\RaggedRight\arraybackslash}X>{\RaggedRight\arraybackslash}X}
\toprule
\textbf{} & \textbf{Planck-only} & \textbf{Combined} & \textbf{Combined (narrow)} \tabularnewline
H\textsubscript{0} (km/\allowbreak{}s/\allowbreak{}Mpc) & 66.89 $\pm$ 0.62 & 67.62 $\pm$ 0.40 & 67.64 $\pm$ 0.40 \tabularnewline
\bottomrule
\end{tabularx}
\end{table}

H\textsubscript{0} is rock-stable across all runs and independent of k\textsubscript{f} prior choice.

\subsubsection{Per-probe breakdown (combined narrow best-fit)}
\begin{table}[H]
\centering
\footnotesize
\begin{tabularx}{\linewidth}{l>{\RaggedRight\arraybackslash}X>{\RaggedRight\arraybackslash}X>{\RaggedRight\arraybackslash}X}
\toprule
\textbf{Probe} & \textbf{$\chi$$^{2}$} & \textbf{N\textsubscript{data}} & \textbf{$\chi$$^{2}$/N} \tabularnewline
Planck plik-lite TTTEEE & 582.74 & 613 & 0.951 \tabularnewline
DESI Y1 BAO & 20.27 & 12 & 1.689 \tabularnewline
Pantheon+ SNe Ia & 2,003.44 & 1,701 & 1.178 \tabularnewline
\textbf{Total} & \textbf{2,606.89} & \textbf{2,326} & \textbf{1.121} \tabularnewline
\bottomrule
\end{tabularx}
\end{table}

\begin{mdframed}[linecolor=accent,linewidth=0.5pt,backgroundcolor=callout-bg,leftline=true,rightline=true,topline=true,bottomline=true,innerleftmargin=8pt,innerrightmargin=8pt,innertopmargin=6pt,innerbottommargin=6pt]
\textbf{Phase 2 conclusion.}
MTDF is fully compatible with Planck 2018 TTTEEE + DESI BAO + Pantheon+ SNe. The predicted stress-field correction (\~{}0.13\% on H(z)) is below current data sensitivity. Full MTDF (\( k_f = 1 \)) and $\Lambda$CDM (\( k_f = 0 \)) are both allowed. Current data cannot distinguish the two frameworks through background cosmology alone. With the physics-motivated narrow prior [0,\,2], the posterior median is \( k_f = 0.93 \), essentially the theoretical prediction. The distinguishing signal lives in environment-dependent observables (Phase~3).
\end{mdframed}

\section{Phase 3: SN $\times$ void environment signal}
\label{sec:sec4}
PASSED: Exact reproduction; bootstrap 95\% CI excludes zeroPhase 3 independently reproduces the SN~Ia~$\times$~void environment analysis reported in MTDF\_\allowbreak{}02 (Companion Part~I). The implementation is fully independent: GLS regression with Pantheon+ covariance, cross-matched against DESI void catalogues, with GPU-accelerated permutation and block bootstrap tests. Runtime: 8.2 minutes (10,000 permutations + 5,000 bootstraps).

\subsection{Environment coefficient $\gamma$\textsubscript{env}}
\begin{table}[H]
\centering
\footnotesize
\begin{tabularx}{\linewidth}{l>{\RaggedRight\arraybackslash}X>{\RaggedRight\arraybackslash}X>{\RaggedRight\arraybackslash}X>{\RaggedRight\arraybackslash}X>{\RaggedRight\arraybackslash}X}
\toprule
\textbf{Catalogue} & \textbf{$\gamma$\textsubscript{env} (mag)} & \textbf{$\Delta$$\chi$$^{2}$} & \textbf{p-value} & \textbf{Target $\Delta$$\chi$$^{2}$} & \textbf{Target p} \tabularnewline
VoidFinder & +0.0030 $\pm$ 0.0018 & 2.92 & 0.088 & 2.91 & 0.088 \tabularnewline
REVOLVER & +0.0047 $\pm$ 0.0023 & 4.25 & 0.039 & 4.25 & 0.039 \tabularnewline
VIDE & +0.0046 $\pm$ 0.0026 & 3.15 & 0.076 & 3.15 & 0.076 \tabularnewline
\bottomrule
\end{tabularx}
\end{table}

All three void catalogues yield positive $\gamma$\textsubscript{env}, matching MTDF\_\allowbreak{}02 values to the last reported decimal place.

\subsection{NGC/\allowbreak{}SGC split}
\begin{table}[H]
\centering
\footnotesize
\begin{tabularx}{\linewidth}{l>{\RaggedRight\arraybackslash}X>{\RaggedRight\arraybackslash}X}
\toprule
\textbf{Footprint} & \textbf{$\gamma$\textsubscript{env} (REVOLVER)} & \textbf{Interpretation} \tabularnewline
NGC & +0.0050 & Strong positive signal \tabularnewline
SGC & +0.0008 & Null, consistent with reduced void coverage \tabularnewline
\bottomrule
\end{tabularx}
\end{table}

The NGC/\allowbreak{}SGC contrast is consistent with footprint-dependent detectability under a finite coherence-length stress field. Dedicated diagnostic tests are planned (see Section 7).

\subsection{Survey fixed effects}
All three catalogues show |$\Delta$$\gamma$\textsubscript{env}| < 0.0005 when survey fixed effects are included. The environment signal is not an artefact of survey-dependent systematics.

\subsection{Leave-one-survey-out (LOSO)}
16 of 16 individual survey removals retain positive $\gamma$\textsubscript{env}. The signal is stable under jackknife.

\subsection{Redshift modulation: the signature result}
\begin{table}[H]
\centering
\footnotesize
\begin{tabularx}{\linewidth}{l>{\RaggedRight\arraybackslash}X>{\RaggedRight\arraybackslash}X>{\RaggedRight\arraybackslash}X>{\RaggedRight\arraybackslash}X}
\toprule
\textbf{Redshift bin} & \textbf{REVOLVER p} & \textbf{VIDE p} & \textbf{VoidFinder p} & \textbf{Interpretation} \tabularnewline
z $\in$ [0.02, 0.04) & \textbf{0.0035} & \textbf{0.0045} & 0.06 & Strong signal \tabularnewline
z > 0.04 & > 0.84 & > 0.84 & > 0.84 & Null \tabularnewline
\bottomrule
\end{tabularx}
\end{table}

\begin{mdframed}[linecolor=accent,linewidth=0.5pt,backgroundcolor=callout-bg,leftline=true,rightline=true,topline=true,bottomline=true,innerleftmargin=8pt,innerrightmargin=8pt,innertopmargin=6pt,innerbottommargin=6pt]
\textbf{Physical interpretation.}
The signal concentrates at low redshift (z~<~0.04) and disappears at higher redshifts. This is a specific, a priori MTDF prediction arising from the finite coherence length of the stress field ($\beta$~=~22.7~Mpc). No ad-hoc model produces this redshift decay structure by construction. $\Lambda$CDM has no mechanism for environment-dependent supernova brightness residuals.
\end{mdframed}

\subsection{Non-parametric confirmation}

\subsubsection{Permutation test (10,000 shuffles)}
\begin{table}[H]
\centering
\footnotesize
\begin{tabularx}{\linewidth}{l>{\RaggedRight\arraybackslash}X}
\toprule
\textbf{Catalogue} & \textbf{p\textsubscript{perm}} \tabularnewline
REVOLVER & 0.025 \tabularnewline
VIDE & 0.046 \tabularnewline
VoidFinder & 0.062 \tabularnewline
\bottomrule
\end{tabularx}
\end{table}

\subsubsection{Block bootstrap (5,000 resamples), 95\% CI}
\begin{table}[H]
\centering
\footnotesize
\begin{tabularx}{\linewidth}{l>{\RaggedRight\arraybackslash}X>{\RaggedRight\arraybackslash}X}
\toprule
\textbf{Catalogue} & \textbf{95\% CI (mag)} & \textbf{Excludes zero?} \tabularnewline
REVOLVER & {}[+0.0016, +0.0076] & Yes \tabularnewline
VIDE & {}[+0.0010, +0.0080] & Yes \tabularnewline
\bottomrule
\end{tabularx}
\end{table}

\subsection{Numerical integrity}
GPU vs CPU crosscheck: maximum difference = 1.86~$\times$~10\textsuperscript{$-$17} (machine epsilon). Results are bit-reproducible.

\begin{mdframed}[linecolor=accent,linewidth=0.5pt,backgroundcolor=callout-bg,leftline=true,rightline=true,topline=true,bottomline=true,innerleftmargin=8pt,innerrightmargin=8pt,innertopmargin=6pt,innerbottommargin=6pt]
\textbf{Phase 3 conclusion.}
The SN~$\times$~void environment signal is independently confirmed. Three void catalogues, 16/\allowbreak{}16 LOSO surveys, permutation tests, and block bootstrap all support a positive $\gamma$\textsubscript{env} with the redshift decay structure predicted by MTDF. Bootstrap 95\% confidence intervals exclude zero. This constitutes the third independent confirmation of the signal (original MTDF\_\allowbreak{}02, independent reproduction, non-parametric validation).
\end{mdframed}

\section{Phase 4: Sensitivity forecasts}
\label{sec:sec5}
COMPLETE: 5$\sigma$ detection at \~{}3,400 SNe (LSST/\allowbreak{}Rubin territory)Phase 4 uses GPU-accelerated Monte Carlo simulations to forecast the sample sizes required for 3$\sigma$ and 5$\sigma$ detection of the environment signal under various systematic scenarios. All six consistency checks against Phase 3 pass (relative differences \~{}10\textsuperscript{$-$13}).

\subsection{Key structural insight}
\begin{mdframed}[linecolor=accent,linewidth=0.5pt,backgroundcolor=callout-bg,leftline=true,rightline=true,topline=true,bottomline=true,innerleftmargin=8pt,innerrightmargin=8pt,innertopmargin=6pt,innerbottommargin=6pt]
\textbf{Rank-1 systematic immunity.}
A global systematic floor of the form \( \sigma_{\rm sys}^2 \times \mathbf{1}\mathbf{1}^T \) has \emph{exactly zero} effect on \( \sigma_{\gamma} \). The intercept absorbs global systematics entirely. Only systematics correlated with \( d_{\rm signed} \) (the void proximity metric) can degrade the environment coefficient. Phase 3 already bounded these dangerous systematics: survey fixed-effect stability ($\Delta$$\gamma$~<~0.0005), LOSO (16/\allowbreak{}16 positive), and redshift modulation (signal at z~<~0.04 only).
\end{mdframed}

\textbf{The regression is statistics-limited, not systematics-limited.}

\subsection{Detection threshold table}
\begin{table}[H]
\centering
\footnotesize
\begin{tabularx}{\linewidth}{l>{\RaggedRight\arraybackslash}X>{\RaggedRight\arraybackslash}X>{\RaggedRight\arraybackslash}X>{\RaggedRight\arraybackslash}X>{\RaggedRight\arraybackslash}X}
\toprule
\textbf{Finder} & \textbf{Scenario} & \textbf{3$\sigma$ 50\%} & \textbf{3$\sigma$ 95\%} & \textbf{5$\sigma$ 50\%} & \textbf{5$\sigma$ 95\%} \tabularnewline
REVOLVER & Baseline & 1,200 & 3,000 & 3,400 & 6,600 \tabularnewline
REVOLVER & Realistic & 1,200 & 2,900 & 3,300 & 6,600 \tabularnewline
REVOLVER & Adversarial & >10k & >10k & >10k & >10k \tabularnewline
REVOLVER & NGC-only & 3,200 & 7,300 & 8,300 & >10k \tabularnewline
VoidFinder & Baseline & 1,700 & 4,500 & 4,900 & 9,100 \tabularnewline
VIDE & Baseline & 1,500 & 4,400 & 4,500 & 8,700 \tabularnewline
\bottomrule
\end{tabularx}
\end{table}

\subsection{Interpretation}
\begin{itemize}
  \item \textbf{Current Pantheon+ (564 SNe):} \~{}2$\sigma$, consistent with Phase 3 ($\Delta$$\chi$$^{2}$~=~4.25 for REVOLVER).
  \item \textbf{Baseline $\approx$ Realistic:} nearly identical thresholds confirm the analysis is statistics-limited, not systematics-limited.
  \item \textbf{Adversarial:} a hypothetical systematic perfectly correlated with \( d_{\rm signed} \) would prevent detection at any sample size, but Phase 3 already rules this out empirically.
  \item \textbf{LSST/\allowbreak{}Rubin (10,000+ low-z SNe):} well into 5$\sigma$ detection territory for all three void catalogues.
\end{itemize}

\begin{mdframed}[linecolor=accent,linewidth=0.5pt,backgroundcolor=callout-bg,leftline=true,rightline=true,topline=true,bottomline=true,innerleftmargin=8pt,innerrightmargin=8pt,innertopmargin=6pt,innerbottommargin=6pt]
\textbf{Merged-sample confirmation:} A merged Pantheon+ / ZTF DR2 / Foundation DR1 sample (4,188 SNe, 3,082 at z < 0.157) confirms that the constant-$\gamma$ global model is null ($\Delta$$\chi$\textsuperscript{2} < 1.4), while a piecewise transition at z $\approx$ 0.04 yields $\Delta$$\chi$\textsuperscript{2} = 36--39 (p < 0.001) across all three void catalogues. SALT2 population controls shift $\gamma$ by < 0.3$\sigma$. The signal is localised to the coherence regime, not a global tilt. See MTDF\_\allowbreak{}02 \S{}2.3 for full merged-sample hardening results.
\end{mdframed}

\begin{mdframed}[linecolor=accent,linewidth=0.5pt,backgroundcolor=callout-bg,leftline=true,rightline=true,topline=true,bottomline=true,innerleftmargin=8pt,innerrightmargin=8pt,innertopmargin=6pt,innerbottommargin=6pt]
\textbf{Phase 4 conclusion.}
Definitive 5$\sigma$ detection of the MTDF environment signal requires approximately 3,400 low-redshift supernovae cross-matched with void catalogues (REVOLVER baseline, 50\% power). This is within the expected yield of LSST/\allowbreak{}Rubin. The analysis is statistics-limited: realistic systematics do not materially degrade detection prospects.
\end{mdframed}

\section{5b. Phase 3b: Asymmetry and detectability diagnostics}
\label{sec:sec5b}
COMPLETE: SGC sensitivity-limited; asymmetry not significant; geometry excludedTo resolve the observed North--South Galactic Cap (NGC/\allowbreak{}SGC) asymmetry in the Phase~3 environment signal, we implemented an expanded diagnostic battery designed to distinguish physical inconsistency from detectability limits imposed by current survey coverage. All five tests ran in production mode (20s total). All diagnostics reuse the exact Phase~3 environment definition (d\textsubscript{signed} array) and regression structure. No void assignments, distances, or environment metrics are recomputed.

\subsection{5b.1 Test A: IPW matched redshift}
After inverse probability weighting to equalise footprint demographics (redshift, host mass proxy, survey origin), no catalogue produces a significant $\Delta$$\gamma$. All p-values range from 0.36 to 0.99. REVOLVER shows the largest residual at +0.0063, still far from significance. Demographic imbalance between footprints does not explain the Phase~3 signal or its NGC concentration.

\subsection{5b.2 Test B: Void--SN joint density mapping}
The SGC is severely void-sparse relative to the NGC across all three catalogues:

\begin{table}[H]
\centering
\footnotesize
\begin{tabularx}{\linewidth}{l>{\RaggedRight\arraybackslash}X>{\RaggedRight\arraybackslash}X>{\RaggedRight\arraybackslash}X}
\toprule
\textbf{Catalogue} & \textbf{NGC voids} & \textbf{SGC voids} & \textbf{Ratio} \tabularnewline
VoidFinder & 264 & 35 & 7.5$\times$ \tabularnewline
REVOLVER & 511 & 82 & 6.2$\times$ \tabularnewline
VIDE & 383 & 58 & 6.6$\times$ \tabularnewline
\bottomrule
\end{tabularx}
\end{table}

SN-per-void ratios are 3--8$\times$ higher in the SGC, indicating that each SGC void is sampled by far fewer supernovae. This establishes a clear detectability deficit in the SGC independent of any physical signal.

\subsection{5b.3 Test C: Wald test and parametric bootstrap}
The Wald test for $\Delta$$\gamma$ = $\gamma$\textsubscript{NGC} $-$ $\gamma$\textsubscript{SGC} returns p-values of 0.49--0.96 across catalogues, indicating no significant asymmetry. The parametric bootstrap (likelihood ratio test under H\textsubscript{0}: $\gamma$\textsubscript{NGC} = $\gamma$\textsubscript{SGC}) returns p = 0.17--0.24. Neither test rejects the null hypothesis that both footprints share the same underlying $\gamma$\textsubscript{env}.

\subsection{5b.4 Test D: Geometry and mode-coupling check}
Isotropic mock catalogues preserving the angular mask and redshift selection function produce $\Delta$$\gamma$ distributions fully consistent with zero. Empirical p-values for the observed $\Delta$$\gamma$ range from 0.28 to 0.90. Survey geometry alone cannot generate the observed hemispherical difference.

\subsection{5b.5 Test E: Null signal injection (SGC recovery)}
An NGC-strength $\gamma$\textsubscript{env} signal was injected into SGC supernova distance moduli using the exact Phase~3 d\textsubscript{signed} values:

\begin{table}[H]
\centering
\footnotesize
\begin{tabularx}{\linewidth}{l>{\RaggedRight\arraybackslash}X>{\RaggedRight\arraybackslash}X>{\RaggedRight\arraybackslash}X}
\toprule
\textbf{Catalogue} & \textbf{Recovery bias} & \textbf{Detection rate at 1$\times$ (2$\sigma$)} & \textbf{Detection rate at 1.5$\times$} \tabularnewline
REVOLVER & <3\% & 44\% & -- \tabularnewline
VIDE & -- & 16\% & -- \tabularnewline
VoidFinder & -- & \~{}0\% & -- \tabularnewline
\bottomrule
\end{tabularx}
\end{table}

VoidFinder's NGC $\gamma$ is near zero, so injection at 1$\times$ is negligible. REVOLVER is the only catalogue with meaningful SGC recovery power.

REVOLVER recovers the injected signal with less than 3\% bias, but achieves only 44\% detection rate at 2$\sigma$ for a 1$\times$ NGC-strength injection. This mechanically demonstrates that the SGC lacks the statistical power to detect the signal even when it is present.

\subsection{5b.6 Interpretation}
\begin{mdframed}[linecolor=accent,linewidth=0.5pt,backgroundcolor=callout-bg,leftline=true,rightline=true,topline=true,bottomline=true,innerleftmargin=8pt,innerrightmargin=8pt,innertopmargin=6pt,innerbottommargin=6pt]
\textbf{Phase 3b conclusion.}
Three findings, each supported by independent diagnostics:
\begin{enumerate}
  \item \textbf{Not significant.} The NGC/\allowbreak{}SGC contrast is not statistically significant (Wald p = 0.49--0.96; bootstrap p = 0.17--0.24). After IPW demographic matching, all $\Delta$$\gamma$ are insignificant (p = 0.36--0.99).
  \item \textbf{Not geometric.} Isotropic mocks with real mask and selection produce null $\Delta$$\gamma$ distributions consistent with zero (p = 0.28--0.90). Survey geometry cannot generate the observed pattern.
  \item \textbf{Sensitivity-limited.} The SGC has 6--8$\times$ fewer voids than the NGC. Null injection tests confirm that an NGC-strength signal is recoverable in SGC-like data only at 44\% power (REVOLVER, 2$\sigma$). The SGC null result reflects insufficient void--SN density, not physical inconsistency with MTDF.
\end{enumerate}
\end{mdframed}

\section{Phase 5: Full Planck plik + class\_\allowbreak{}mtdf}
\label{sec:sec6}
PASSED: $\Delta$$\chi$$^{2}$ = +0.63; k\textsubscript{f} = 0.028 $\approx$ 0.

Phase 5 is the Priority~1 hard falsifier: running the full Planck 2018 plik TTTEEE + lowl TT + lowl EE + lensing likelihood using the \texttt{class\_\allowbreak{}mtdf} modified Boltzmann solver, which computes modified perturbation equations, transfer functions, ISW contributions, and lensing potentials directly. This is not a perturbative approximation; the full MTDF physics propagates through every C\textsubscript{$\ell$} multipole. Minimization was performed with cobaya's BOBYQA algorithm (best of 4 parallel runs), floating 6 cosmological parameters, 1 MTDF parameter (k\textsubscript{f}), and 21 Planck nuisance parameters.

\begin{mdframed}[linecolor=accent,linewidth=0.5pt,backgroundcolor=callout-bg,leftline=true,rightline=true,topline=true,bottomline=true,innerleftmargin=8pt,innerrightmargin=8pt,innertopmargin=6pt,innerbottommargin=6pt]
\textbf{Independence caveat.} Phase 5 is not independent in the same sense as Phases 1, 3, 4, and 6, because it uses the existing \texttt{class\_\allowbreak{}mtdf} solver. Its role is to test that solver against the full Planck likelihood, not to provide a completely independent reimplementation of the Boltzmann equations.
\end{mdframed}

\begin{mdframed}[linecolor=accent,linewidth=0.5pt,backgroundcolor=callout-bg,leftline=true,rightline=true,topline=true,bottomline=true,innerleftmargin=8pt,innerrightmargin=8pt,innertopmargin=6pt,innerbottommargin=6pt]
\textbf{Apples-to-apples guarantee.}
Both $\Lambda$CDM and MTDF runs use identical: likelihood set (plik TTTEEE + lowl TT + lowl EE + lensing), data bins, masks, nuisance parameter list (21 sampled), nuisance priors (Planck-shipped), cosmological priors, and minimizer settings (BOBYQA best\_\allowbreak{}of=4). The only difference is that the MTDF run adds k\textsubscript{f} as a free parameter and uses \texttt{class\_\allowbreak{}mtdf} instead of standard CLASS.
\end{mdframed}

\subsection{$\chi$$^{2}$ breakdown by likelihood component}
\begin{table}[H]
\centering
\footnotesize
\begin{tabularx}{\linewidth}{l>{\RaggedRight\arraybackslash}X>{\RaggedRight\arraybackslash}X>{\RaggedRight\arraybackslash}X}
\toprule
\textbf{Component} & \textbf{$\Lambda$CDM} & \textbf{MTDF} & \textbf{$\Delta$(MTDF $-$ $\Lambda$CDM)} \tabularnewline
lowl TT & 23.25 & 25.90 & +2.65 \tabularnewline
lowl EE & 395.69 & 396.08 & +0.39 \tabularnewline
plik TTTEEE & 2,345.41 & 2,341.74 & $-$3.67 \tabularnewline
Lensing & 8.85 & 10.10 & +1.25 \tabularnewline
\textbf{Total} & \textbf{2,773.20} & \textbf{2,773.82} & \textbf{+0.63} \tabularnewline
\bottomrule
\end{tabularx}
\end{table}

$\Lambda$CDM total $\chi$$^{2}$ = 2,773.20 is consistent with the published Planck 2018 best-fit (\~{}2,780). Pipeline is correctly calibrated.

The high-$\ell$ plik TTTEEE likelihood actually \emph{improves} by $-$3.67 for MTDF, offset by a degradation in lowl TT of +2.65. The lensing likelihood shifts by +1.25. The net effect is a total $\Delta$$\chi$$^{2}$ of +0.63, which is statistically negligible.

\subsection{Model selection}
\begin{table}[H]
\centering
\footnotesize
\begin{tabularx}{\linewidth}{l>{\RaggedRight\arraybackslash}X>{\RaggedRight\arraybackslash}X>{\RaggedRight\arraybackslash}X>{\RaggedRight\arraybackslash}X}
\toprule
\textbf{Metric} & \textbf{$\Lambda$CDM} & \textbf{MTDF} & \textbf{$\Delta$} & \textbf{Verdict} \tabularnewline
$\Delta$$\chi$$^{2}$ & -- & -- & +0.63 & Essentially tied \tabularnewline
AIC (k = 27 vs 28) & 2,827.20 & 2,829.82 & +2.63 & Borderline $\Lambda$CDM \tabularnewline
BIC (N $\approx$ 1,676) & 2,973.65 & 2,981.70 & +8.05 & $\Lambda$CDM preferred \tabularnewline
\bottomrule
\end{tabularx}
\end{table}

BIC penalises the extra parameter more heavily. This is expected and correct: Planck alone does not need k\textsubscript{f}. The MTDF signal lives in late-universe probes.

\subsection{Best-fit cosmological parameters}
\begin{table}[H]
\centering
\footnotesize
\begin{tabularx}{\linewidth}{l>{\RaggedRight\arraybackslash}X>{\RaggedRight\arraybackslash}X>{\RaggedRight\arraybackslash}X}
\toprule
\textbf{Parameter} & \textbf{$\Lambda$CDM} & \textbf{MTDF} & \textbf{$\Delta$} \tabularnewline
H\textsubscript{0} (km/\allowbreak{}s/\allowbreak{}Mpc) & 67.257 & 67.282 & +0.025 \tabularnewline
n\textsubscript{s} & 0.9650 & 0.9650 & $-$0.0001 \tabularnewline
$\omega$\textsubscript{b} & 0.02234 & 0.02235 & +0.00002 \tabularnewline
$\omega$\textsubscript{cdm} & 0.12022 & 0.12018 & $-$0.00004 \tabularnewline
$\tau$\textsubscript{reio} & 0.0507 & 0.0538 & +0.003 \tabularnewline
$\sigma$\textsubscript{8} & 0.8088 & 0.7946 & $-$0.014 \tabularnewline
k\textsubscript{f} & -- & 0.028 & -- \tabularnewline
\bottomrule
\end{tabularx}
\end{table}

\begin{mdframed}[linecolor=accent,linewidth=0.5pt,backgroundcolor=callout-bg,leftline=true,rightline=true,topline=true,bottomline=true,innerleftmargin=8pt,innerrightmargin=8pt,innertopmargin=6pt,innerbottommargin=6pt]
\textbf{Note on $\sigma$\textsubscript{8}.}
At the minimizer best-fit, MTDF returns $\sigma$\textsubscript{8} = 0.795 vs $\Lambda$CDM's 0.809 (1.7\% shift). The full MCMC posteriors confirm a larger separation: $\sigma$\textsubscript{8} = 0.790 $\pm$ 0.006 (MTDF) vs 0.810 $\pm$ 0.006 ($\Lambda$CDM), a 2.4\% shift relative to the $\Lambda$CDM posterior mean (\~{}2.4$\sigma$ combining 1$\sigma$ uncertainties in quadrature, Gaussian approximation). This is in the direction needed to ease the S\textsubscript{8} tension between Planck and weak lensing surveys (KiDS, DES). The shift was not imposed; it emerged naturally from the sampler.
\end{mdframed}

\subsection{Interpretation}
The optimizer found k\textsubscript{f} = 0.028 $\approx$ 0, confirming that Planck CMB data alone cannot activate the MTDF stress-field correction. This is the expected result: MTDF predicts that its corrections are concentrated at low redshift via the finite coherence length, making them invisible to the CMB. The framework passes the most demanding CMB test available.

This result was obtained by numerical minimization (BOBYQA), not from an MCMC posterior mean. A full MCMC posterior sampling run is planned as a follow-up to provide proper credible intervals and parameter degeneracy maps.

\begin{mdframed}[linecolor=accent,linewidth=0.5pt,backgroundcolor=callout-bg,leftline=true,rightline=true,topline=true,bottomline=true,innerleftmargin=8pt,innerrightmargin=8pt,innertopmargin=6pt,innerbottommargin=6pt]
\textbf{Phase 5 conclusion (minimization).}
MTDF passes the full Planck 2018 plik TTTEEE + lowl + lensing test with $\Delta$$\chi$$^{2}$ = +0.63 using the \texttt{class\_\allowbreak{}mtdf} Boltzmann solver with full modified perturbation equations. This is the Priority~1 hard falsifier identified in the test programme. The result confirms Phase~2's perturbative finding ($\Delta$$\chi$$^{2}$~<~1) and shows that this MTDF implementation is statistically indistinguishable from $\Lambda$CDM at the CMB within the precision and assumptions of this likelihood comparison, even when the full physics of modified transfer functions, ISW, and lensing are included.
\end{mdframed}

\subsection{Planck degeneracy structure and posterior stability (MCMC converged)}
Full MCMC posterior sampling has converged for both $\Lambda$CDM and MTDF chains using cobaya's strict two-stage convergence criteria. Both chains reached Gelman--Rubin R$-$1~<~0.02 with credible interval convergence R$-$1\textsubscript{cl}~<~0.07.

\subsubsection{Convergence statistics}
\begin{table}[H]
\centering
\footnotesize
\begin{tabularx}{\linewidth}{l>{\RaggedRight\arraybackslash}X>{\RaggedRight\arraybackslash}X}
\toprule
\textbf{Metric} & \textbf{$\Lambda$CDM} & \textbf{MTDF} \tabularnewline
R$-$1 (final) & 0.019 & 0.019 \tabularnewline
R$-$1\textsubscript{cl} (final) & 0.064 & 0.057 \tabularnewline
Accepted samples & 26,400 & 27,160 \tabularnewline
ESS (H\textsubscript{0}) & \~{}1,920 & \~{}1,450 \tabularnewline
ESS ($\sigma$\textsubscript{8}) & \~{}2,020 & \~{}1,600 \tabularnewline
ESS (k\textsubscript{f}) & -- & \~{}1,470 \tabularnewline
\bottomrule
\end{tabularx}
\end{table}

\subsubsection{Posterior parameter comparison (68\% CI)}
\begin{table}[H]
\centering
\footnotesize
\begin{tabularx}{\linewidth}{l>{\RaggedRight\arraybackslash}X>{\RaggedRight\arraybackslash}X>{\RaggedRight\arraybackslash}X}
\toprule
\textbf{Parameter} & \textbf{$\Lambda$CDM} & \textbf{MTDF} & \textbf{$\Delta$} \tabularnewline
H\textsubscript{0} (km/\allowbreak{}s/\allowbreak{}Mpc) & 67.38 $\pm$ 0.54 & 67.83 $\pm$ 0.54 & +0.45 \tabularnewline
$\sigma$\textsubscript{8} & 0.810 $\pm$ 0.006 & 0.790 $\pm$ 0.006 & $-$0.020 \tabularnewline
S\textsubscript{8} & 0.830 & 0.802 & $-$0.028 \tabularnewline
$\Omega$\textsubscript{m} & 0.315 $\pm$ 0.007 & 0.309 $\pm$ 0.007 & $-$0.006 \tabularnewline
n\textsubscript{s} & 0.965 $\pm$ 0.004 & 0.968 $\pm$ 0.004 & +0.004 \tabularnewline
$\tau$\textsubscript{reio} & 0.054 $\pm$ 0.007 & 0.052 $\pm$ 0.007 & $-$0.002 \tabularnewline
k\textsubscript{f} & -- & 0.50 $\pm$ 0.36 & -- \tabularnewline
\bottomrule
\end{tabularx}
\end{table}

\subsubsection{k\textsubscript{f} posterior}
The marginalised k\textsubscript{f} posterior is broad and approximately Gaussian: k\textsubscript{f}~=~0.50~$\pm$~0.36, with 95\%~CI~=~[0.025,~1.34]. k\textsubscript{f}~=~1 (full MTDF) lies within the 95\% credible interval. k\textsubscript{f}~=~0 ($\Lambda$CDM) is not strongly excluded and sits near the lower boundary, which is shaped by the non-negativity prior (k\textsubscript{f}~$\ge$~0). Planck does not constrain k\textsubscript{f} tightly, confirming that MTDF is a late-time modification invisible to the CMB.

\subsubsection{k\textsubscript{f} identifiability diagnostics}
Identifiability diagnostics show that k\textsubscript{f} does not participate in the dominant Planck degeneracy blocks and is not a reparameterisation of any single nuisance parameter. All Planck nuisance correlations satisfy |r|~<~0.10, and k\textsubscript{f} shows negligible correlation with the amplitude--optical-depth combination A\textsubscript{s}$\cdot$e\textsuperscript{$-$2$\tau$} (r~$\approx$~0.09). Two-dimensional posteriors are close to round against H\textsubscript{0}, log\,A and $\tau$, with only mild covariance against $\sigma$\textsubscript{8}. Nuisance parameters alone explain 13.7\% of k\textsubscript{f} variance (R$^{2}$~=~0.137); early-time primordial parameters explain 27.8\% (R$^{2}$~=~0.278), consistent with a parameter weakly constrained by the CMB. The marginal posterior is data-informed, peaking near k\textsubscript{f}~$\approx$~0.15--0.20 under a uniform prior on [0,\,5], with k\textsubscript{f}~=~1 inside the 95\% credible interval and k\textsubscript{f}~=~0 lying near the lower boundary under the non-negativity prior.

k\textsubscript{f} shows a moderate covariance with $\omega$\textsubscript{b} (marginal r~=~$-$0.30; stable partial r~$\approx$~$-$0.34 conditioning on n\textsubscript{s} and $\omega$\textsubscript{cdm}). Full conditioning yields an inflated partial correlation (+0.97) because the conditioning matrix is ill-conditioned ($\kappa$~=~24,207; multiple predictors with VIF~>~100), so we treat that as a multicollinearity artefact rather than a physical one-to-one degeneracy. Detailed diagnostics including the progressive conditioning table, residual-method cross-check, and VIF report are archived in \texttt{output/\allowbreak{}phase5/\allowbreak{}robustness/\allowbreak{}kf\_\allowbreak{}identifiability/\allowbreak{}}.

\subsubsection{Best-fit robustness across Planck subsets}
Best-fit BOBYQA comparisons (best-of-4, likelihood-only $\chi$$^{2}$) across Planck subsets show $\Delta$$\chi$$^{2}$ values close to zero: +0.63 (full baseline), +0.41 (no lensing), and $-$0.65 (TT only). The TT-only subset marginally favours MTDF at best fit, while all subsets yield $\Delta$AIC~<~3, indicating no meaningful preference for either model. Corresponding $\Delta$AIC values ($\Delta$k~=~1) are +2.63, +2.41, and +1.35. These results are consistent with MTDF neither degrading nor improving the Planck fit across likelihood configurations.

\subsubsection{Phase 5 robustness checks (Tests 2--4)}
We performed three complementary robustness diagnostics to assess whether the inferred MTDF coupling parameter k\textsubscript{f} and the associated shifts in derived parameters are artefacts of a particular Planck likelihood component, nuisance freedom, or prior choice. All robustness artefacts are packaged under \texttt{output/\allowbreak{}phase5/\allowbreak{}robustness/\allowbreak{}} with manifests and canonical JSON summaries.

\textbf{Test 2: Leave-one-likelihood-out (Planck subsets).}
Repeating the MTDF inference with lensing removed yields results consistent with the baseline posterior: k\textsubscript{f}~=~0.527~$\pm$~0.377 versus baseline 0.495~$\pm$~0.360, with $\sigma$\textsubscript{8} and H\textsubscript{0} shifts remaining small in standard-deviation units ($\Delta$k\textsubscript{f}~$\approx$~+0.09$\sigma$, $\Delta$$\sigma$\textsubscript{8}~$\approx$~+0.64$\sigma$). A TT-only run exhibits the expected weakening of constraints and broader posteriors, with k\textsubscript{f}~=~0.938~$\pm$~0.670 and $\sigma$\textsubscript{8} uncertainty significantly enlarged; this reflects reduced constraining power in TT alone rather than a change in physical behaviour. Best-fit BOBYQA comparisons using likelihood-only $\chi$$^{2}$ remain near parity across subsets: baseline $\Delta$$\chi$$^{2}$~=~+0.63, no lensing $\Delta$$\chi$$^{2}$~=~+0.41, and TT only $\Delta$$\chi$$^{2}$~=~$-$0.65 (definitions and values given in the Test~2 artefacts).

\textbf{Test 4: Prior range sensitivity.}
Widening the k\textsubscript{f} prior from [0,\,5] to [0,\,10] leaves the posterior essentially unchanged: k\textsubscript{f}~=~0.483~$\pm$~0.361 (wide) versus 0.495~$\pm$~0.360 (baseline), with nearly identical 95\% credible intervals ([0.023,\,1.341] versus [0.025,\,1.342]) and a negligible mean shift ($-$0.03$\sigma$). This indicates that the Phase~5 k\textsubscript{f} inference is data-informed and not driven by the prior bound.

\textbf{Summary.}
Taken together with the identifiability diagnostics in Test~3, these results show that the Phase~5 posterior behaviour is stable under removal of lensing, behaves as expected under the information loss of TT only, and is insensitive to widening the k\textsubscript{f} prior. The robustness suite therefore supports interpreting the observed $\sigma$\textsubscript{8} suppression and modest H\textsubscript{0} shift as stable consequences of the MTDF coupling within the Planck likelihood framework, rather than numerical artefacts or prior-driven effects.

\begin{mdframed}[linecolor=accent,linewidth=0.5pt,backgroundcolor=callout-bg,leftline=true,rightline=true,topline=true,bottomline=true,innerleftmargin=8pt,innerrightmargin=8pt,innertopmargin=6pt,innerbottommargin=6pt]
\textbf{Phase 5 MCMC converged.}
Both chains reached R$-$1~<~0.02 (threshold: 0.02). The $\sigma$\textsubscript{8} shift ($-$2.4\% relative to the $\Lambda$CDM posterior mean, \~{}2.4$\sigma$ in Gaussian approximation) and directional H\textsubscript{0} shift (+0.45~km~s\textsuperscript{$-$1}~Mpc\textsuperscript{$-$1}, \~{}0.6$\sigma$) first observed in Phase~2 (plik-lite) are confirmed at full Planck precision with identical nuisance freedom. k\textsubscript{f}~=~1 (full MTDF) lies within the 95\% credible interval; k\textsubscript{f}~=~0 ($\Lambda$CDM) is not strongly excluded and sits near the lower boundary. Triangle plots and k\textsubscript{f} posterior are archived in \texttt{output/\allowbreak{}phase5/\allowbreak{}}.
\end{mdframed}

\section{6b. Phase 6: Cross-probe discriminator tests}
\label{sec:sec6b}
COMPLETE: Two discriminators (A, E), one coherence check (C2), three upper limits (B, B2, D), one internal consistency (C), two forecasts (C3, C4), one resolved null (C5)  -  10 tests totalPhase~6 extends validation beyond CMB compatibility toward tests that can \emph{distinguish} MTDF from $\Lambda$CDM using independent probes. Ten tests are packaged under \texttt{output/\allowbreak{}phase6/\allowbreak{}}, targeting weak lensing, CMB lensing, derived parameter consistency, cross-probe S\textsubscript{8} coherence, growth--environment coupling, BAO residuals, cluster mass ratios, and the SN~$\times$~void redshift transition. All code was written from scratch for Phase~6.

Tests~A and E positively support MTDF-style predictions (redshift transition, void-specific CMB lensing signal). Test~C2 demonstrates cross-probe S\textsubscript{8} coherence (combined weak lensing tension reduced from 3.5$\sigma$ to 1.9$\sigma$). Tests~B, B2, and D are systematics-clean nulls with explicit upper limits. Test~C confirms internal consistency. Tests~C3 and C4 quantify sensitivity boundaries. Test~C5 resolves the cluster mass ratio ($\eta$~=~1). Together they demonstrate: \emph{we did not cheat, the pipelines are sound, no hidden contradictions exist, and the boundaries of current discriminating power are mapped.}

\subsection{Summary table}
\begin{table}[H]
\centering
\footnotesize
\begin{tabularx}{\linewidth}{l>{\RaggedRight\arraybackslash}X>{\RaggedRight\arraybackslash}X>{\RaggedRight\arraybackslash}X>{\RaggedRight\arraybackslash}X}
\toprule
\textbf{Test} & \textbf{Observable} & \textbf{Primary statistic} & \textbf{Result} & \textbf{Implication} \tabularnewline
\midrule
\textbf{A} & SN~$\times$~void $\gamma$\textsubscript{env} vs z\textsubscript{cut} & GLS z-score peak over z\textsubscript{cut} scan & Peak at z\textsubscript{cut}=0.030: 3.62$\sigma$ (global p~<~0.005) & Supports MTDF transition at z~<~0.04 \tabularnewline
\textbf{B} & KiDS-1000 $\gamma$\textsubscript{t} around void centres & $\Delta$$\gamma$\textsubscript{t} over R=5--20 Mpc/\allowbreak{}h & $-$3.65e$-$5~$\pm$~4.68e$-$5 (0.8$\sigma$); 95\%: |$\Delta$$\gamma$\textsubscript{t}|~<~1.2e$-$4 & Null  -  bounded, not yet discriminating \tabularnewline
\textbf{B2} & KiDS-1000 trough vs ridge $\gamma$\textsubscript{t} & $\Delta$$\gamma$\textsubscript{t,split} (ridge$-$trough) & +2.56e$-$5~$\pm$~5.46e$-$5 (0.5$\sigma$); 95\%: |$\Delta$$\gamma$\textsubscript{t,split}|~<~1.3e$-$4 & Null  -  bounded, not yet discriminating \tabularnewline
\textbf{C} & Derived parameter consistency & Max shift vs $\Lambda$CDM; S\textsubscript{8} tension reduction & All within 2.4$\sigma$; S\textsubscript{8}: DES 2.6$\rightarrow$1.3$\sigma$, KiDS 2.8$\rightarrow$1.6$\sigma$ & Internally consistent; eases S\textsubscript{8} tension \tabularnewline
\textbf{D} & Planck PR4 $\kappa$ $\times$ DESI BGS voids & Per-void compensated $\kappa$\textsubscript{comp} & $-$7.47e$-$4~$\pm$~1.03e$-$3 (0.7$\sigma$, 557 voids) & Null at low z  -  expected; pipeline validated \tabularnewline
\textbf{E} & Planck PR4 $\kappa$ $\times$ BOSS DR12 voids & Conservative $\kappa$\textsubscript{comp} + detection-optimised CTH/\allowbreak{}MF & Comp: 0.2$\sigma$ null; CTH: 1.5$\sigma$ (p\textsubscript{rand}=0.03); MF: 1.9$\sigma$ (p\textsubscript{rand}=0.02) & Pipeline benchmark validated; void-specific signal recovered \tabularnewline
\textbf{C2} & S\textsubscript{8} cross-probe coherence & Phase~5 S\textsubscript{8} vs published WL surveys & Combined WL tension: 3.47$\sigma$~$\rightarrow$~1.89$\sigma$ (46\% reduction) & MTDF shifts S\textsubscript{8} in the right direction by the right amount \tabularnewline
\textbf{C3} & BAO residual structure & $\Delta$W RMS (wiggle difference MTDF$-$$\Lambda$CDM) & $\Delta$W RMS~=~0.27\%, max~=~0.51\%; sound horizon shift +0.11\% & Not detectable by any foreseeable survey (best SNR~=~0.92) \tabularnewline
\textbf{C4} & f$\sigma$\textsubscript{8} $\times$ environment & RSD + void f$\sigma$\textsubscript{8} vs $\mu$(a) prediction & $\Delta$$\chi$$^{2}$~=~+0.49 (9 DOF); max discrimination 0.22$\sigma$ & Sensitivity-limited; needs \~{}10$\times$ improvement (DESI/\allowbreak{}Euclid) \tabularnewline
\textbf{C5} & Cluster M\textsubscript{dyn}/M\textsubscript{lens} & Analytical $\mu$(a) prediction + numerical $\eta$(k,z) verification & $\eta$~=~1.000000 at all cluster scales; Case~B confirmed: M\textsubscript{dyn}/M\textsubscript{lens}~=~1 & Resolved: cluster mass ratio is not a discriminator for MTDF \tabularnewline
\bottomrule
\end{tabularx}
\end{table}

\subsection{6b.1 Test A: Redshift transition scan}
\begin{itemize}
  \item Scans z\textsubscript{cut} $\in$ [0.025,~0.100] using GLS and Spearman metrics on the Phase~3 SN~$\times$~void environment signal.
  \item Peak at z\textsubscript{cut}~=~0.030 with GLS z-score 3.62$\sigma$; global scan-corrected permutation p~<~0.005 (0 of 200 exceedances).
  \item Spearman rank test confirms independently (3.1$\sigma$).
  \item Interpretation: the environment signal is localised to z~<~0.04, consistent with the Phase~3 redshift modulation and the MTDF coherence-length prediction.
  \item Artefacts: output/\allowbreak{}phase6/\allowbreak{}testA\_\allowbreak{}redshift\_\allowbreak{}transition/\allowbreak{}
\end{itemize}

\subsection{6b.2 Test B: Weak lensing $\times$ environment (void centres)}
\begin{itemize}
  \item KiDS-1000 shear stacked around DESI REVOLVER void centres and matched random positions; $\gamma$\textsubscript{x} (cross-shear) gate required to pass before reporting $\gamma$\textsubscript{t}.
  \item $\Delta$$\gamma$\textsubscript{t}~=~$-$3.65$\times$10\textsuperscript{$-$5}~$\pm$~4.68$\times$10\textsuperscript{$-$5} over R~=~5--20~Mpc/\allowbreak{}h, consistent with zero.
  \item $\gamma$\textsubscript{x} gate: pass (p~=~0.234).
  \item Explicit two-sided 95\% upper limit: |$\Delta$$\gamma$\textsubscript{t}|~<~1.20$\times$10\textsuperscript{$-$4}.
  \item Without a forward MTDF lensing prediction for this selection, the result is best framed as ``not yet discriminating, but now bounded and reproducible.''
  \item Artefacts: output/\allowbreak{}phase6/\allowbreak{}testB\_\allowbreak{}wl\_\allowbreak{}environment/\allowbreak{}
\end{itemize}

\subsection{6b.3 Test B2: Trough and ridge lensing (projected density split)}
\begin{itemize}
  \item Projected density field built with mask-aware normalisation and full KiDS-1000 footprint (North + South). Both trough and ridge $\gamma$\textsubscript{x} gates pass.
  \item $\Delta$$\gamma$\textsubscript{t,split} (ridge$-$trough)~=~+2.56$\times$10\textsuperscript{$-$5}~$\pm$~5.46$\times$10\textsuperscript{$-$5} (0.5$\sigma$), consistent with zero.
  \item Explicit two-sided 95\% upper limit: |$\Delta$$\gamma$\textsubscript{t,split}|~<~1.28$\times$10\textsuperscript{$-$4}.
  \item Projected P20/\allowbreak{}P80 density contrasts ($\delta$~$\sim$~$\pm$0.15) are too mild for detectable differential lensing at KiDS sensitivity. Spectroscopic 3D density or LSST/\allowbreak{}DES area would be needed for significance.
  \item Artefacts: output/\allowbreak{}phase6/\allowbreak{}testB2\_\allowbreak{}trough\_\allowbreak{}ridge/\allowbreak{}
\end{itemize}

\subsection{6b.4 Test C: Derived parameter consistency}
\begin{itemize}
  \item Cross-checks derived parameters from the Phase~5 MCMC posterior against external weak lensing constraints.
  \item No derived parameter deviates beyond 2.4$\sigma$ from $\Lambda$CDM.
  \item S\textsubscript{8} tension reduction: DES~Y3 from 2.63$\sigma$ to 1.28$\sigma$; KiDS-1000 from 2.77$\sigma$ to 1.58$\sigma$.
  \item MTDF's Planck posterior is internally consistent; no hidden contradictions introduced by \texttt{class\_\allowbreak{}mtdf}.
  \item Artefacts: output/\allowbreak{}phase6/\allowbreak{}testC\_\allowbreak{}derived\_\allowbreak{}consistency/\allowbreak{}
\end{itemize}

\subsection{6b.5 Test D: CMB lensing $\times$ DESI BGS voids (low redshift)}
\begin{itemize}
  \item Planck PR4 MV convergence ($\kappa$) stacked on DESIVAST BGS voids (z~<~0.24) with mask-weighted stacking. Primary statistic is a per-void compensated filter: $\kappa$(R/\allowbreak{}R\textsubscript{v}<0.5)~$-$~$\kappa$(4<R/\allowbreak{}R\textsubscript{v}<5), designed to remove large-scale mode contamination.
  \item $\kappa$\textsubscript{comp}~=~$-$7.47$\times$10\textsuperscript{$-$4}~$\pm$~1.03$\times$10\textsuperscript{$-$3} (0.7$\sigma$, 557 voids), consistent with zero.
  \item RA-scramble null: p~=~0.425 (compensated). Random-sky null: p~=~0.080.
  \item Robustness: stable across outer annulus choices (all <1.1$\sigma$), low-l removal (l$\ge$20, l$\ge$30), and void finder.
  \item Null at z~<~0.24 is expected: the CMB lensing kernel peaks at z~\~{}~1--2, and the BGS void count (557 usable) limits statistical power.
  \item Diagnosed failure modes from v1: hard mask quality cut rejected 73\% of voids; low-l CMB modes created a positive $\kappa$ pedestal. Both fixed in v2 via per-void compensated filter and mask-weighted stacking.
  \item Artefacts: output/\allowbreak{}phase6/\allowbreak{}testD\_\allowbreak{}cmb\_\allowbreak{}lensing/\allowbreak{}
\end{itemize}

\subsection{6b.6 Test E: CMB lensing $\times$ BOSS DR12 voids (benchmark)}
\begin{itemize}
  \item Benchmark run on the Mao+2017 ZOBOV catalogue (1,228 BOSS DR12 CMASS+LOWZ voids, z~=~0.2--0.7) to validate the CMB lensing stacking pipeline against Cai+2017's published 3.2$\sigma$ detection on the same catalogue.
  \item Five statistics computed side by side:
    
\textbf{Per-void compensated} (systematics-clean primary): 0.2$\sigma$, null by design.
\textbf{AP disc} (contamination monitor): 4.7$\sigma$, but RA-scramble p~=~0.945 and collapses to 0.3$\sigma$ under l$\ge$20 cut  -  entirely pedestal.
\textbf{CTH} (Cai-class compensated top-hat): 1.5$\sigma$ (p\textsubscript{rand}~=~0.03), increases to 2.0$\sigma$ under l$\ge$20 cut.
\textbf{Matched filter} (template + jackknife covariance): 1.9$\sigma$ (p\textsubscript{rand}~=~0.02), increases to 2.2$\sigma$ under l$\ge$20 cut.
\textbf{Close comp} (alternative aperture): 0.5$\sigma$, consistent with zero.
  \item Detection-optimised statistics (CTH and MF) show the expected hierarchy and respond correctly to pedestal removal and footprint-preserving null tests. The gap from \~{}1.5--2.2$\sigma$ to Cai+2017's 3.2$\sigma$ reflects template shape (step vs theory-derived) and different Planck products (PR4 vs PR2015).
  \item Artefacts: output/\allowbreak{}phase6/\allowbreak{}testE\_\allowbreak{}boss\_\allowbreak{}benchmark/\allowbreak{}
\end{itemize}

\subsection{6b.7 Test C2: S\textsubscript{8} cross-probe coherence}
\begin{itemize}
  \item Compares Phase~5 MCMC S\textsubscript{8}~=~$\sigma$\textsubscript{8}$\sqrt{\,}$($\Omega$\textsubscript{m}/0.3) against published weak lensing constraints from KiDS-1000, DES~Y3, and HSC~Y3.
  \item $\Lambda$CDM: S\textsubscript{8}~=~0.830~$\pm$~0.013. MTDF: S\textsubscript{8}~=~0.802~$\pm$~0.012.
  \item Tension reduction: KiDS-1000 2.67$\sigma$~$\rightarrow$~1.63$\sigma$ (39\%); DES~Y3 2.63$\sigma$~$\rightarrow$~1.28$\sigma$ (51\%); combined WL 3.47$\sigma$~$\rightarrow$~1.89$\sigma$ (46\%).
  \item Caveat: published WL S\textsubscript{8} values assume $\Lambda$CDM growth and lensing kernel. A proper comparison requires WL likelihood within the MTDF framework (deferred to future work).
  \item Artefacts: output/\allowbreak{}phase6/\allowbreak{}testC2\_\allowbreak{}s8\_\allowbreak{}coherence/\allowbreak{}
\end{itemize}

\subsection{6b.8 Test C3: BAO residual structure}
\begin{itemize}
  \item Extracts BAO wiggles W(k)~=~P(k)/P\textsubscript{nw}(k)~$-$~1 from \texttt{class\_\allowbreak{}mtdf} power spectra using an Eisenstein~\&~Hu (1998) no-wiggle reference, and computes the residual $\Delta$W~=~W\textsubscript{MTDF}~$-$~W\textsubscript{$\Lambda$CDM}.
  \item Sound horizon shift: +0.106\% (147.148~$\rightarrow$~147.303~Mpc). Wiggle difference: $\Delta$W RMS~=~0.27\%, max~=~0.51\% at k~=~0.12~h/\allowbreak{}Mpc. No detectable peak phase shift at grid resolution.
  \item Best projected SNR: 0.92 (DESI + Euclid combined). The 0.1\% sound horizon shift is absorbed into standard cosmological parameter degeneracies.
  \item Conclusion: BAO residual structure is not a viable MTDF--$\Lambda$CDM discriminator with any foreseeable survey.
  \item Artefacts: output/\allowbreak{}phase6/\allowbreak{}testC3\_\allowbreak{}bao\_\allowbreak{}residuals/\allowbreak{}
\end{itemize}

\subsection{6b.9 Test C4: f$\sigma$\textsubscript{8} $\times$ environment}
\begin{itemize}
  \item Compares MTDF and $\Lambda$CDM f$\sigma$\textsubscript{8}(z) predictions against a compilation of 8 RSD surveys plus void-specific growth measurements from Hamaus+2020.
  \item $\chi$$^{2}$: $\Lambda$CDM~=~7.61, MTDF~=~8.10 ($\Delta$$\chi$$^{2}$~=~+0.49, 9~DOF). Maximum per-point discrimination: 0.22$\sigma$.
  \item The homogeneous $\mu$(a) modification produces a \~{}1.5\% signal, which is 5--10$\times$ smaller than current RSD measurement errors (0.03--0.10 per z-bin).
  \item Future requirement: 2$\sigma$ discrimination needs $\sigma$(f$\sigma$\textsubscript{8})~\~{}~0.003--0.004 per z-bin, achievable with DESI/\allowbreak{}Euclid-era RSD or environment-resolved estimators.
  \item Artefacts: output/\allowbreak{}phase6/\allowbreak{}testC4\_\allowbreak{}fsigma8\_\allowbreak{}environment/\allowbreak{}
\end{itemize}

\subsection{6b.10 Test C5/\allowbreak{}C5b: Cluster M\textsubscript{dyn}/M\textsubscript{lens}  -  Resolved}
\begin{itemize}
  \item C5 computed the predicted ratio M\textsubscript{dyn}/M\textsubscript{lens} under two bracketing coupling cases:
    
\textbf{Case~A} ($\mu$-only dynamics, $\Sigma$~=~1): M\textsubscript{dyn}/M\textsubscript{lens}~=~1.053 at z~=~0 (+5.3\% offset).
\textbf{Case~B} (equal coupling, $\Sigma$~=~$\mu$): M\textsubscript{dyn}/M\textsubscript{lens}~=~1.000 (no offset).
  \item \textbf{C5b resolved the key unknown} by extracting $\eta$(k,z)~=~$\Phi$(k,z)/$\Psi$(k,z) directly from \texttt{class\_\allowbreak{}mtdf} Newtonian-gauge transfer functions across k~=~0.005--2~h/\allowbreak{}Mpc and z~=~0--10.
  \item Result: \textbf{$\eta$~=~1.000000} at all cluster-relevant scales (z~$\le$~0.8). Max~|$\Delta$$\eta$| between MTDF and $\Lambda$CDM:~0.000000. Tiny deviations (|$\eta$$-$1|~\~{}~10\textsuperscript{$-$5}) appear only at z~$\ge$~5 from standard radiation anisotropic stress, present equally in both models.
  \item Physical reason: \texttt{class\_\allowbreak{}mtdf} modifies only the Poisson equation source ($\delta$$\rho$~$\rightarrow$~$\mu$$\cdot$$\delta$$\rho$) and the CDM Euler equation. The $\Phi$$-$$\Psi$ anisotropy equation is unchanged from GR. Both potentials are enhanced equally $\rightarrow$ $\Sigma$~=~$\mu$ $\rightarrow$ \textbf{Case~B confirmed}.
  \item Implication: M\textsubscript{dyn}/M\textsubscript{lens}~=~1.000 under MTDF. The cluster mass ratio channel is \textbf{not a discriminator}. Clusters can still test MTDF through growth-dependent observables (e.g.~cluster counts), because $\mu$~>~1 modifies the growth rate.
  \item Artefacts: output/\allowbreak{}phase6/\allowbreak{}testC5\_\allowbreak{}cluster\_\allowbreak{}mass/\allowbreak{}, output/\allowbreak{}phase6/\allowbreak{}testC5b\_\allowbreak{}gravitational\_\allowbreak{}slip/\allowbreak{}
\end{itemize}

\begin{mdframed}[linecolor=accent,linewidth=0.5pt,backgroundcolor=callout-bg,leftline=true,rightline=true,topline=true,bottomline=true,innerleftmargin=8pt,innerrightmargin=8pt,innertopmargin=6pt,innerbottommargin=6pt]
\textbf{Phase 6 conclusion (10 tests, all complete).}
The ten Phase~6 tests fall into five categories. \emph{Discriminators} (6A, 6E): the redshift transition is detected at 3.6$\sigma$; the BOSS CMB lensing benchmark recovers void-specific structure at 1.5--1.9$\sigma$. \emph{Coherence} (6C, 6C2): derived parameters are self-consistent (no shift > 2.4$\sigma$); S\textsubscript{8} tension with weak lensing surveys drops from 3.5$\sigma$ to 1.9$\sigma$. \emph{Upper limits} (6B, 6B2, 6D): systematics-clean nulls with explicit 95\% bounds on void lensing and CMB lensing signals. \emph{Forecasts} (6C3, 6C4): BAO wiggles and f$\sigma$\textsubscript{8} quantify what future surveys need to reach. \emph{Resolved} (6C5): cluster M\textsubscript{dyn}/M\textsubscript{lens} is null ($\eta$~=~1). No test contradicts MTDF. The S\textsubscript{8} coherence result (6C2) is the strongest positive outcome: MTDF moves S\textsubscript{8} toward weak lensing values without tuning.
\end{mdframed}

\section{Synthesis and framework narrative}
\label{sec:sec7}

\subsection{Summary of results (17/\allowbreak{}17 complete)}
\begingroup\footnotesize
\begin{longtable}{>{\RaggedRight\arraybackslash}p{0.192\linewidth}>{\RaggedRight\arraybackslash}p{0.192\linewidth}>{\RaggedRight\arraybackslash}p{0.192\linewidth}>{\RaggedRight\arraybackslash}p{0.192\linewidth}>{\RaggedRight\arraybackslash}p{0.192\linewidth}}
\toprule
\endfirsthead
\toprule\endhead
\bottomrule\endfoot
\textbf{Phase} & \textbf{Type} & \textbf{Status} & \textbf{Key result} & \textbf{Artefacts} \tabularnewline
Phase 1 & Discriminator & PASSED & V18 strict totals reproduced to $\delta$~=~0.000049 & output/\allowbreak{}phase1/\allowbreak{} \tabularnewline
Phase 2 & Discriminator & PASSED & CMB plik-lite compatible ($\Delta$$\chi$$^{2}$~<~1); k\textsubscript{f}~=~1 within 95\% in all MCMC runs & output/\allowbreak{}phase2/\allowbreak{} \tabularnewline
Phase 3 & Discriminator & PASSED & SN~$\times$~void signal reproduced exactly; bootstrap excludes zero & output/\allowbreak{}phase3/\allowbreak{} \tabularnewline
Phase 3b & Upper limit & COMPLETE & NGC/\allowbreak{}SGC asymmetry: not significant, not geometric, sensitivity-limited (5 diagnostics) & output/\allowbreak{}phase3b/\allowbreak{} \tabularnewline
Phase 4 & Forecast & COMPLETE & 5$\sigma$ at \~{}3,400 SNe (LSST/\allowbreak{}Rubin); statistics-limited & output/\allowbreak{}phase4/\allowbreak{} \tabularnewline
Phase 5 (min.) & Discriminator & PASSED & Full Planck plik + class\_\allowbreak{}mtdf: $\Delta$$\chi$$^{2}$~=~+0.63. Priority~1 hard falsifier cleared. & \multirow{2}{*}{output/\allowbreak{}phase5/\allowbreak{}} \tabularnewline
Phase 5 (MCMC) & Discriminator & PASSED & Converged (R$-$1~=~0.019). $\sigma$\textsubscript{8}~=~0.790 (\~{}2.4$\sigma$ below $\Lambda$CDM). k\textsubscript{f}~=~0.50~$\pm$~0.36; both k\textsubscript{f}~=~0 and 1 within 95\% CI. &  \tabularnewline
Phase 6A & Discriminator & PASSED & Redshift transition scan: signal peaks at z~<~0.03 (3.6$\sigma$ GLS, global p~<~0.005). Spearman confirms (3.1$\sigma$). & \multirow{10}{*}{output/\allowbreak{}phase6/\allowbreak{}} \tabularnewline
Phase 6B & Upper limit & COMPLETE & Weak lensing $\times$ void environment (KiDS-1000): $\Delta$$\gamma$\textsubscript{t}~=~$-$3.65e$-$5~$\pm$~4.68e$-$5 (0.8$\sigma$). $\gamma$\textsubscript{x} gate pass. 95\% UL: |$\Delta$$\gamma$\textsubscript{t}|~<~1.2e$-$4. &  \tabularnewline
Phase 6B2 & Upper limit & COMPLETE & Trough/\allowbreak{}ridge lensing (KiDS-1000): $\Delta$$\gamma$\textsubscript{t,split}~=~+2.56e$-$5~$\pm$~5.46e$-$5 (0.5$\sigma$). Both $\gamma$\textsubscript{x} gates pass. 95\% UL: |$\Delta$$\gamma$\textsubscript{t,split}|~<~1.3e$-$4. &  \tabularnewline
Phase 6C & Coherence & PASSED & Derived parameter consistency: no shifts >~2.4$\sigma$. S\textsubscript{8} tension reduces: DES~Y3 2.6$\rightarrow$1.3$\sigma$; KiDS-1000 2.8$\rightarrow$1.6$\sigma$. &  \tabularnewline
Phase 6C2 & Coherence & COMPLETE & S\textsubscript{8} cross-probe coherence: Phase~5 posterior S\textsubscript{8} vs WL surveys. Combined tension 3.5$\sigma$~$\rightarrow$~1.9$\sigma$ (46\% reduction). Summary-level; full likelihood deferred. &  \tabularnewline
Phase 6C3 & Forecast & COMPLETE & BAO residual structure: $\Delta$W RMS~=~0.27\%, sound horizon shift +0.11\%. Not detectable by any current or projected survey (best SNR~=~0.92). &  \tabularnewline
Phase 6C4 & Forecast & COMPLETE & f$\sigma$\textsubscript{8}~$\times$~environment: RSD + void growth vs $\mu$(a) prediction. Max 0.22$\sigma$; needs DESI/\allowbreak{}Euclid-era precision (\~{}10$\times$ improvement). &  \tabularnewline
Phase 6C5/\allowbreak{}C5b & Closed null & RESOLVED & Cluster M\textsubscript{dyn}/M\textsubscript{lens}: $\eta$~=~1 confirmed numerically (C5b) $\Rightarrow$ $\Sigma$~=~$\mu$, so M\textsubscript{dyn}/M\textsubscript{lens}~=~1 exactly. Mass ratio offset is zero in the current MTDF perturbation model. &  \tabularnewline
Phase 6D & Upper limit & COMPLETE & CMB lensing $\times$ DESI BGS voids: $\kappa$\textsubscript{comp}~=~$-$7.47e$-$4~$\pm$~1.03e$-$3 (0.7$\sigma$, 557 voids). RA-scramble clean. Null expected at z~<~0.24. &  \tabularnewline
Phase 6E & Discriminator & COMPLETE & CMB lensing $\times$ BOSS DR12 benchmark: CTH 1.5$\sigma$ (p\textsubscript{rand}=0.03), MF 1.9$\sigma$ (p\textsubscript{rand}=0.02). Pipeline validated against Cai+2017. &  \tabularnewline
\end{longtable}
\endgroup

\subsection{Claim status taxonomy}
Each MTDF claim or test outcome is assigned one of four status labels. This taxonomy is intended to provide referees and readers with an unambiguous assessment of where the framework stands.

\begin{table}[H]
\centering
\footnotesize
\begin{tabularx}{\linewidth}{l>{\RaggedRight\arraybackslash}X}
\toprule
\textbf{Label} & \textbf{Definition} \tabularnewline
\checkmark{} Validated & Independently reproduced with quantitative agreement. Robust to systematics tests. \tabularnewline
$\blacksquare$ Constrained & Consistent with data but not uniquely favoured. Current data cannot distinguish from $\Lambda$CDM. \tabularnewline
$\blacktriangle$ Hypothesis & Theoretically motivated but not yet tested against data, or tested only indirectly. \tabularnewline
$\times$ Falsifier & Identified test that could refute the framework. Not yet executed. \tabularnewline
\bottomrule
\end{tabularx}
\end{table}

\subsubsection{Status by claim}
\begingroup\footnotesize
\begin{longtable}{>{\RaggedRight\arraybackslash}p{0.320\linewidth}>{\RaggedRight\arraybackslash}p{0.320\linewidth}>{\RaggedRight\arraybackslash}p{0.320\linewidth}}
\toprule
\endfirsthead
\toprule\endhead
\bottomrule\endfoot
\textbf{Claim / element} & \textbf{Status} & \textbf{Evidence} \tabularnewline
V18 strict $\chi$$^{2}$ totals & \checkmark{} Validated & Phase 1: independent reproduction to $\delta$ = 0.000049 \tabularnewline
Pantheon+ SN likelihood & \checkmark{} Validated & Phase 1: $\chi$$^{2}$ = 2012.72 matched exactly \tabularnewline
DESI Y1 BAO fit & \checkmark{} Validated & Phase 1: $\chi$$^{2}$ = 15.39 matched exactly \tabularnewline
Cosmic chronometers H(z) & \checkmark{} Validated & Phase 1: $\chi$$^{2}$ = 6.98 matched exactly \tabularnewline
DR16 f$\sigma$\textsubscript{8} growth & \checkmark{} Validated & Phase 1: $\chi$$^{2}$ = 2.00 matched exactly \tabularnewline
CMB compatibility (plik-lite) & $\blacksquare$ Constrained & Phase 2: $\Delta$$\chi$$^{2}$ < 1; k\textsubscript{f} = 1 within 95\% CI. Perturbative correction layer, not full Boltzmann. \tabularnewline
CMB compatibility (full plik + class\_\allowbreak{}mtdf) & \checkmark{} Validated & Phase 5: full plik TTTEEE + lowl + lensing via class\_\allowbreak{}mtdf Boltzmann solver. $\Delta$$\chi$$^{2}$ = +0.63 (BOBYQA best-fit, 4 runs). k\textsubscript{f} = 0.028 $\approx$ 0. Apples-to-apples with $\Lambda$CDM (same likelihoods, nuisances, priors, minimizer). \tabularnewline
SN $\times$ void environment signal ($\gamma$\textsubscript{env}) & \checkmark{} Validated & Phase 3: exact reproduction, 3 void finders, bootstrap excludes zero \tabularnewline
Redshift decay structure (z < 0.04) & \checkmark{} Validated & Phase 3: a priori prediction confirmed (p = 0.0035 REVOLVER) \tabularnewline
Survey systematics stability & \checkmark{} Validated & Phase 3: FE $\Delta$$\gamma$ < 0.0005; LOSO 16/\allowbreak{}16 positive \tabularnewline
NGC/\allowbreak{}SGC asymmetry interpretation & \checkmark{} Validated & Phase 3b: asymmetry not significant (Wald p = 0.49--0.96), not geometric (mock p = 0.28--0.90), SGC sensitivity-limited (6--8$\times$ fewer voids, 44\% recovery power). Five independent diagnostics confirm footprint-dependent detectability. \tabularnewline
5$\sigma$ detection forecast (3,400 SNe) & \checkmark{} Validated & Phase 4: Monte Carlo with consistency checks passing at 10\textsuperscript{$-$13} \tabularnewline
Statistics-limited (not systematics-limited) & \checkmark{} Validated & Phase 4: baseline $\approx$ realistic; rank-1 systematic immunity demonstrated \tabularnewline
Hubble tension resolution (void depletion) & $\blacksquare$ Constrained & Two independent equations yield H\textsubscript{local} consistent with SH0ES. Mechanism independent of CMB; not yet tested with dedicated void-SN joint analysis. \tabularnewline
Photon--stress coupling ($\kappa$ = f\textsubscript{kick}/3 $\approx$ 0.00102) & $\blacktriangle$ Hypothesis & Derived from fundamental parameters (angular averaging of strain tensor). Below Planck sensitivity. Paper MTDF\_\allowbreak{}04 frames as speculative. \tabularnewline
Rotation curve phenomenology & $\blacksquare$ Constrained & Included in V18 scalars. Not independently reproduced in this session. \tabularnewline
Redshift transition structure (z\textsubscript{cut} scan) & \checkmark{} Validated & Phase 6A: GLS z-score 3.62$\sigma$ at z\textsubscript{cut}=0.030; global scan-corrected p~<~0.005. Spearman confirms (3.1$\sigma$). \tabularnewline
Weak lensing $\times$ environment (KiDS) & $\blacksquare$ Constrained & Phase 6B/\allowbreak{}B2: systematics-clean nulls with explicit 95\% upper limits. Bounded but not yet discriminating at KiDS sensitivity. \tabularnewline
Derived parameter consistency & \checkmark{} Validated & Phase 6C: all within 2.4$\sigma$. S\textsubscript{8} tension eased (DES 2.6$\rightarrow$1.3$\sigma$; KiDS 2.8$\rightarrow$1.6$\sigma$). \tabularnewline
CMB lensing $\times$ voids & $\blacksquare$ Constrained & Phase 6D/\allowbreak{}E: low-z null (expected), BOSS benchmark recovers void-specific structure. Pipeline validated. \tabularnewline
S\textsubscript{8} cross-probe coherence & \checkmark{} Validated & Phase 6C2: S\textsubscript{8} from Phase~5 posterior vs WL surveys. Combined tension 3.5$\sigma$~$\rightarrow$~1.9$\sigma$. Summary-level comparison; full likelihood deferred. \tabularnewline
f$\sigma$\textsubscript{8} $\times$ environment & $\blacksquare$ Constrained & Phase 6C4: $\mu$(a) prediction vs RSD + void growth. Max 0.22$\sigma$ discrimination; current data sensitivity-limited (\~{}10$\times$ improvement needed). \tabularnewline
BAO residual structure & $\blacksquare$ Constrained & Phase 6C3: $\Delta$W RMS = 0.27\%. Sound horizon shift +0.11\% absorbed into parameter degeneracies. Not detectable by any foreseeable survey. \tabularnewline
Cluster M\textsubscript{dyn}/M\textsubscript{lens} & $\blacksquare$ Constrained & Phase 6C5/\allowbreak{}C5b: $\eta$~=~1.000000 confirmed numerically from class\_\allowbreak{}mtdf transfer functions. Case~B applies: M\textsubscript{dyn}/M\textsubscript{lens}~=~1. Mass ratio channel closed as discriminator; growth-dependent cluster tests remain open. \tabularnewline
Weak lensing P(k) cross-correlation & $\times$ Falsifier & Not yet executed. Requires full P(k) from class\_\allowbreak{}mtdf at non-linear scales. \tabularnewline
Merging cluster offsets (Bullet Cluster) & $\times$ Falsifier & Not yet executed. Requires N-body simulations. \tabularnewline
\end{longtable}
\endgroup

\subsection{Framework narrative}
\begin{mdframed}[linecolor=accent,linewidth=0.5pt,backgroundcolor=callout-bg,leftline=true,rightline=true,topline=true,bottomline=true,innerleftmargin=8pt,innerrightmargin=8pt,innertopmargin=6pt,innerbottommargin=6pt]
MTDF is fully CMB-compatible: Planck TTTEEE cannot distinguish MTDF from $\Lambda$CDM ($\Delta$$\chi$$^{2}$~<~1). MTDF deviates from $\Lambda$CDM at late times through environment-dependent dynamics. The Hubble tension is resolved by environmental mapping (void stress depletion), not by early-universe modification. The SN~$\times$~void signal is the primary empirical discriminator. The stress-field coherence length concentrates observable effects at low redshift, consistent with the CMB being unaffected.
\end{mdframed}

\subsection{What survives intact}
\begin{enumerate}
  \item \textbf{Phase 1 strict totals}: independently reproduced to four decimal places.
  \item \textbf{SN $\times$ void environment signal}: three void finders, 16/\allowbreak{}16 LOSO, redshift decay matching a priori prediction.
  \item \textbf{Hubble tension resolution}: entirely late-time, via void stress depletion:
    
Theoretical: \( H_{\rm local}/H_0 = 1.084 \;\Rightarrow\; 67.4 \times 1.084 = 73.1 \) (SH0ES: 73.04~$\pm$~1.04)
MCMC-derived: \( H_{0,\rm global} = 68.56 \) + void enhancement 5.81 = \( H_{0,\rm local} = 74.37 \)
  \item \textbf{Late-time phenomenology}: w(z), BAO distances, growth rate, rotation curves.
  \item \textbf{Framework stability}: MTDF recovers $\Lambda$CDM at early times, deviates at late times.
  \item \textbf{Cross-probe integrity (Phase 6, 10 tests)}: redshift transition confirmed (3.6$\sigma$); weak lensing and CMB lensing pipelines pass strict systematics controls; derived parameters self-consistent; S\textsubscript{8} tension reduced from 3.5$\sigma$ to 1.9$\sigma$ (combined WL); f$\sigma$\textsubscript{8}, BAO residuals, and cluster mass predictions computed  -  all sensitivity-limited with current data, with quantified thresholds for future surveys.
\end{enumerate}

\subsection{Joint implications for late-time cosmological tensions}
While MTDF is not tuned to resolve any single cosmological tension, the validated framework exhibits simultaneous shifts in multiple late-time observables. In particular, the Phase~5 best-fit favours a modest reduction in $\sigma$\textsubscript{8} relative to $\Lambda$CDM best-fit values (0.795 vs 0.809), consistent with the direction preferred by weak-lensing surveys (KiDS, DES).

Independently, the environment-dependent dynamics tested in Phases~3 and 4 naturally modify local expansion observables without compromising global CMB consistency. Together, these results indicate that MTDF addresses multiple late-time phenomenological tensions within a single, unified dynamical framework, without introducing additional free parameters beyond k\textsubscript{f}.

\subsection{Paper implications}
\begin{table}[H]
\centering
\footnotesize
\begin{tabularx}{\linewidth}{l>{\RaggedRight\arraybackslash}X}
\toprule
\textbf{Paper} & \textbf{Impact} \tabularnewline
MTDF\_\allowbreak{}01 (Master) & Phase 1 strengthens all claims. No changes needed. \tabularnewline
MTDF\_\allowbreak{}02 (Part I, Environmental) & Signal independently confirmed (Phase 3). Redshift transition localised to z~<~0.04 (Phase 6A, 3.6$\sigma$). Framework's strongest empirical evidence. \tabularnewline
MTDF\_\allowbreak{}04 (Part III, Photon coupling) & $\kappa$ $\approx$ f\textsubscript{kick}/3 is an observational anchor structurally related to the early field energy fraction and below Planck sensitivity. Paper frames this as speculative, which is appropriate. \tabularnewline
MTDF\_\allowbreak{}05 (Part IV, Cosmological) & Phase 5 full plik total $\Delta$$\chi$$^{2}$~=~+0.63; Phase 2 plik-lite precursor: +0.81; CMB+BAO summary run: +1.25 (narrower data). Hubble tension mechanism fully intact. Priority 1 hard falsifier cleared. \tabularnewline
MTDF\_\allowbreak{}07 (This paper) & Phase 6 adds ten cross-probe tests (6A--6E, 6C2--6C5). Title updated to Phases 1--6. All validation claims current as of February 2026. \tabularnewline
\bottomrule
\end{tabularx}
\end{table}

\section{Open questions and next steps}
\label{sec:sec8}

\subsection{NGC/\allowbreak{}SGC asymmetry diagnostics (Phase 3b, complete)}
Phase 3b (Section 5b) implemented the full asymmetry diagnostic battery and is now complete. All five tests confirm that the NGC/\allowbreak{}SGC contrast is not statistically significant, not caused by survey geometry, and attributable to SGC void-catalogue sparsity. The NGC/\allowbreak{}SGC claim status in the taxonomy (Section 7.2) has been upgraded from $\blacktriangle$~Hypothesis to \checkmark{}~Validated.

\subsection{Falsifier and discriminator tests}
\begin{table}[H]
\centering
\footnotesize
\begin{tabularx}{\linewidth}{l>{\RaggedRight\arraybackslash}X>{\RaggedRight\arraybackslash}X>{\RaggedRight\arraybackslash}X>{\RaggedRight\arraybackslash}X}
\toprule
\textbf{\#} & \textbf{Test} & \textbf{Type} & \textbf{Status} & \textbf{Timescale} \tabularnewline
1 & Full Planck plik + class\_\allowbreak{}mtdf & Discriminator & \textbf{PASSED} (Phase 5, $\Delta$$\chi$$^{2}$ = +0.63) & Complete \tabularnewline
2 & NGC/\allowbreak{}SGC diagnostic battery & Upper limit & \textbf{COMPLETE} (sensitivity-limited, not geometric) & Complete \tabularnewline
3 & Weak lensing $\times$ environment (KiDS) & Upper limit & \textbf{COMPLETE} (Phase 6B/\allowbreak{}B2: systematics-clean null) & Complete \tabularnewline
4 & CMB lensing $\times$ voids & Upper limit & \textbf{COMPLETE} (Phase 6D/\allowbreak{}E: pipeline validated) & Complete \tabularnewline
5 & S\textsubscript{8} cross-probe coherence & Coherence & \textbf{COMPLETE} (Phase 6C2: 3.5$\sigma$~$\rightarrow$~1.9$\sigma$) & Complete \tabularnewline
6 & f$\sigma$\textsubscript{8} $\times$ environment & Forecast & \textbf{COMPLETE} (Phase 6C4: 0.22$\sigma$, sensitivity-limited) & Complete \tabularnewline
7 & BAO residual structure & Forecast & \textbf{COMPLETE} (Phase 6C3: $\Delta$W = 0.27\%, undetectable) & Complete \tabularnewline
8 & Cluster M\textsubscript{dyn}/M\textsubscript{lens} & Resolved & \textbf{RESOLVED} (Phase 6C5/\allowbreak{}C5b: $\eta$~=~1, Case~B, M\textsubscript{dyn}/M\textsubscript{lens}~=~1) & Complete \tabularnewline
9 & Weak lensing P(k) (non-linear) & Discriminator & Requires full P(k) from class\_\allowbreak{}mtdf & Months \tabularnewline
10 & Merging cluster offsets & Discriminator & Requires N-body simulations & Institutional \tabularnewline
\bottomrule
\end{tabularx}
\end{table}

\begin{mdframed}[linecolor=accent,linewidth=0.5pt,backgroundcolor=callout-bg,leftline=true,rightline=true,topline=true,bottomline=true,innerleftmargin=8pt,innerrightmargin=8pt,innertopmargin=6pt,innerbottommargin=6pt]
\textbf{Phase 2 scope note, resolved by Phase 5.}
Phase 2 used a perturbative correction layer on a $\Lambda$CDM emulator, not the full modified Boltzmann equations. Phase~5 subsequently ran \texttt{class\_\allowbreak{}mtdf} through the full Planck plik likelihood with all nuisance parameters, confirming $\Delta$$\chi$$^{2}$~=~+0.63. The perturbative approximation in Phase 2 was validated by the full calculation in Phase 5.
\end{mdframed}

\section{Appendix: Technical notes}
\label{sec:sec9}

\subsection{A.1 \,k\textsubscript{f} normalisation conventions}
Three different \( k_f \) conventions were identified across the MTDF codebase during Phase 2. Understanding these is essential for interpreting historical results.

\begin{table}[H]
\centering
\footnotesize
\begin{tabularx}{\linewidth}{l>{\RaggedRight\arraybackslash}X>{\RaggedRight\arraybackslash}X>{\RaggedRight\arraybackslash}X}
\toprule
\textbf{Convention} & \textbf{"Full MTDF" k\textsubscript{f}} & \textbf{f\textsubscript{kick}} & \textbf{Source} \tabularnewline
A (old CLASS, pre-correction) & 0.10 & 0.003303 & November 2025 cobaya runs \tabularnewline
B (corrected CLASS) & 1.0 & 0.003303 & Theory prediction \tabularnewline
C (Phase 2 code, initial, corrected) & 0.001326 & 0.001326 & CosmoPower correction layer (bug) \tabularnewline
\bottomrule
\end{tabularx}
\end{table}

Convention C initially used \( f_{\rm kick} = \alpha\kappa = 0.001326 \) instead of the correct \( f_{\rm kick} = \lambda_{\rm MTDF}/24 = 0.003303 \). This factor-of-2.49 error caused a 3$\times$ overcorrection in the angular scale ratio per unit \( k_f \), leading to a spurious \( H_0 \)--\( k_f \) degeneracy in the first MCMC run (best-fit \( H_0 = 55.6 \), which is unphysical). After correction, all runs show \( H_0 \approx 67 \) and no degeneracies.

\subsection{A.2 \,V18 parameter reference}
\begin{table}[H]
\centering
\footnotesize
\begin{tabularx}{\linewidth}{l>{\RaggedRight\arraybackslash}X>{\RaggedRight\arraybackslash}X>{\RaggedRight\arraybackslash}X}
\toprule
\textbf{Parameter} & \textbf{Value} & \textbf{Uncertainty} & \textbf{Role} \tabularnewline
$\alpha$ & 1.30 & 0.26 & Stress-field coupling (void dynamics: KBC 2013 + Hoscheit 2018; counter: Kenworthy 2019) \tabularnewline
$\beta$ & 22.7 Mpc & 0.09$\times$$10^{23}$ m & Coherence length (void radius distribution: Pan+2012; Hamaus+2014) \tabularnewline
$\tau$ & 13.0 Gyr & 0.2 & Field age \tabularnewline
$\beta$\textsubscript{eos} & 0.573 & 0.012 & Equation-of-state parameter \tabularnewline
E & 9.1$\times$10\textsuperscript{$-$10} Pa & 0.4$\times$10\textsuperscript{$-$10} & Field energy scale \tabularnewline
$\delta$\textsubscript{bf} & 0.150 & 0.030 & Observational anchor: cluster baryon fraction at M\textsubscript{500} $\sim$ 10\textsuperscript{14} M\textsubscript{$\odot$} (Giodini+2009; Gonzalez+2013; see MTDF\_\allowbreak{}01 \S{}5.1) \tabularnewline
z\textsubscript{t} & 0.74 & 0.10 & Transition redshift \tabularnewline
$\kappa$ & 0.00102 & 0.00003 & Photon--stress coupling (anchor: $\kappa$ $\approx$ f\textsubscript{kick}/3; below FIRAS sensitivity) \tabularnewline
H\textsubscript{0} & 70.0 km/\allowbreak{}s/\allowbreak{}Mpc & 1.5 & Hubble constant (Freedman 2021 TRGB: 69.8 $\pm$ 0.6) \tabularnewline
f\textsubscript{kick} & 0.003303 & -- & Derived: (1$-$$\beta$\textsubscript{eos})$^{2}$/((1+$\alpha$)$\times$24) \tabularnewline
\bottomrule
\end{tabularx}
\end{table}

\subsection{A.3 \,Historical context: November 2025 runs}
Previous cobaya + class\_\allowbreak{}mtdf runs (November 2025) provide useful reference points. A single-point evaluation at the Planck $\Lambda$CDM best-fit gave $\chi$$^{2}$~=~614.23 (consistent with full plik TT+TE+EE). An MCMC with Planck TT-only (16 walkers, 200 steps) found \( k_f \approx 0.03 \) in the old normalisation (Convention A), with $\chi$$^{2}$ indistinguishable from $\Lambda$CDM. These early results are consistent with the current Phase 2 conclusion that Planck cannot constrain the EFE amplitude.

\section{Data and code availability}
\label{sec:data-availability}
The complete MTDF reproducibility package, including the independent GPU
pipeline used in this paper, the modified CLASS solver (\texttt{class\_\allowbreak{}mtdf}),
the cobaya YAML configurations for Phase~5, the Phase~2 CosmoPower emulator
cache, the V18 strict validation workbook, and the gravity-sector artefact
manifest, is archived on Zenodo
(concept DOI \href{https://doi.org/10.5281/zenodo.19741058}{10.5281/\allowbreak{}zenodo.\allowbreak{}19741058};
this release, \texttt{v1.1.2}, DOI
\href{https://doi.org/10.5281/zenodo.19743916}{10.5281/\allowbreak{}zenodo.\allowbreak{}19743916})
and mirrored on GitHub at
\href{https://github.com/IMesche/MTDF}{github.com/\allowbreak{}IMesche/\allowbreak{}MTDF}.
The \texttt{v1.1.2} Git tag corresponds to the archived Zenodo release. All
numerical results reported here, including the Phase~1 reproduction, the
Phase~2 plik-lite MCMC, the Phase~3 SN~$\times$~void bootstrap, the Phase~5
full Planck minimisation and converged MCMC, and the Phase~6 cross-probe
discriminators, can be regenerated from the scripts under
\texttt{gpu\_\allowbreak{}validation/} following the procedure in the
top-level \texttt{README.md}.

\section{References}
\begin{thebibliography}{19}
\bibitem{ref1} Aghanim, N. et al. (Planck Collaboration) 2020, "Planck 2018 results. V. Power spectra and likelihoods", \emph{A\&A}, 641, A5. DOI: \href{https://doi.org/10.1051/0004-6361/201936386}{10.1051/\allowbreak{}0004-6361/\allowbreak{}201936386}
\bibitem{ref2} Aghanim, N. et al. (Planck Collaboration) 2020, "Planck 2018 results. VI. Cosmological parameters", \emph{A\&A}, 641, A6. DOI: \href{https://doi.org/10.1051/0004-6361/201833910}{10.1051/\allowbreak{}0004-6361/\allowbreak{}201833910}
\bibitem{ref3} DESI Collaboration 2025, "DESI 2024 VI: Cosmological Constraints from BAO", \emph{JCAP}, 2025, 021. DOI: \href{https://doi.org/10.1088/1475-7516/2025/02/021}{10.1088/\allowbreak{}1475-7516/\allowbreak{}2025/\allowbreak{}02/\allowbreak{}021}
\bibitem{ref4} Brout, D. et al. 2022, "The Pantheon+ Analysis: Cosmological Constraints", \emph{ApJ}, 938, 110. DOI: \href{https://doi.org/10.3847/1538-4357/ac8e04}{10.3847/\allowbreak{}1538-4357/\allowbreak{}ac8e04}
\bibitem{ref5} Melia, F. \& Maier, R. S. 2013, "Cosmic chronometers in the R\_\allowbreak{}h=ct universe", \emph{MNRAS}, 432, 2669. DOI: \href{https://doi.org/10.1093/mnras/stt1082}{10.1093/\allowbreak{}mnras/\allowbreak{}stt1082}
\bibitem{ref6} Alam, S. et al. (eBOSS Collaboration) 2021, "Completed SDSS-IV eBOSS: Cosmological implications", \emph{Phys. Rev. D}, 103, 083533. DOI: \href{https://doi.org/10.1103/PhysRevD.103.083533}{10.1103/\allowbreak{}PhysRevD.103.083533}
\bibitem{ref7} Spurio Mancini, A. et al. 2022, "CosmoPower: emulating cosmological power spectra for accelerated Bayesian inference", \emph{MNRAS}, 511, 1771. DOI: \href{https://doi.org/10.1093/mnras/stac064}{10.1093/\allowbreak{}mnras/\allowbreak{}stac064}
\bibitem{ref8} Torrado, J. \& Lewis, A. 2021, "Cobaya: Code for Bayesian Analysis of hierarchical physical models", \emph{JCAP}, 2021, 057. DOI: \href{https://doi.org/10.1088/1475-7516/2021/05/057}{10.1088/\allowbreak{}1475-7516/\allowbreak{}2021/\allowbreak{}05/\allowbreak{}057}
\bibitem{ref9} Foreman-Mackey, D. et al. 2013, "emcee: The MCMC Hammer", \emph{PASP}, 125, 306. DOI: \href{https://doi.org/10.1086/670067}{10.1086/\allowbreak{}670067}
\bibitem{ref10} Blas, D., Lesgourgues, J. \& Tram, T. 2011, "The Cosmic Linear Anisotropy Solving System (CLASS). Part II", \emph{JCAP}, 2011, 034. DOI: \href{https://doi.org/10.1088/1475-7516/2011/07/034}{10.1088/\allowbreak{}1475-7516/\allowbreak{}2011/\allowbreak{}07/\allowbreak{}034}
\bibitem{ref11} Asgari, M. et al. (KiDS Collaboration) 2021, "KiDS-1000 Cosmology: Cosmic shear constraints", \emph{A\&A}, 645, A104. DOI: \href{https://doi.org/10.1051/0004-6361/202039070}{10.1051/\allowbreak{}0004-6361/\allowbreak{}202039070}
\bibitem{ref12} Abbott, T. M. C. et al. (DES Collaboration) 2022, "Dark Energy Survey Year 3 results: Cosmological constraints", \emph{Phys. Rev. D}, 105, 023520. DOI: \href{https://doi.org/10.1103/PhysRevD.105.023520}{10.1103/\allowbreak{}PhysRevD.105.023520}
\bibitem{ref13} Dalal, R. et al. (HSC Collaboration) 2023, "Hyper Suprime-Cam Year 3 Results: Cosmology from Cosmic Shear", \emph{Phys. Rev. D}, 108, 123519. DOI: \href{https://doi.org/10.1103/PhysRevD.108.123519}{10.1103/\allowbreak{}PhysRevD.108.123519}
\bibitem{ref14} Keenan, R. C., Barger, A. J. \& Cowie, L. L. 2013, "Evidence for a \~{}300 Megaparsec Scale Under-density", \emph{ApJ}, 775, 62. DOI: \href{https://doi.org/10.1088/0004-637X/775/1/62}{10.1088/\allowbreak{}0004-637X/\allowbreak{}775/\allowbreak{}1/\allowbreak{}62}
\bibitem{ref15} Pan, D. C. et al. 2012, "Cosmic Voids in SDSS DR7", \emph{MNRAS}, 421, 926. DOI: \href{https://doi.org/10.1111/j.1365-2966.2011.20197.x}{10.1111/\allowbreak{}j.1365-2966.2011.20197.x}
\bibitem{ref16} Freedman, W. L. 2021, "Measurements of the Hubble Constant: Tensions in Perspective", \emph{ApJ}, 919, 16. DOI: \href{https://doi.org/10.3847/1538-4357/ac0e95}{10.3847/\allowbreak{}1538-4357/\allowbreak{}ac0e95}
\bibitem{ref17} McGaugh, S. S., Lelli, F. \& Schombert, J. M. 2016, "Radial Acceleration Relation in Rotationally Supported Galaxies", \emph{PRL}, 117, 201101. DOI: \href{https://doi.org/10.1103/PhysRevLett.117.201101}{10.1103/\allowbreak{}PhysRevLett.117.201101}
\bibitem{ref18} Lelli, F. et al. 2017, "One Law to Rule Them All: The Radial Acceleration Relation", \emph{ApJ}, 836, 152. DOI: \href{https://doi.org/10.3847/1538-4357/836/2/152}{10.3847/\allowbreak{}1538-4357/\allowbreak{}836/\allowbreak{}2/\allowbreak{}152}
\bibitem{ref19} Riess, A. G. et al. 2022, "A Comprehensive Measurement of the Local Value of the Hubble Constant", \emph{ApJL}, 934, L7. DOI: \href{https://doi.org/10.3847/2041-8213/ac5c5b}{10.3847/\allowbreak{}2041-8213/\allowbreak{}ac5c5b}
\end{thebibliography}


\end{document}
